% ===========================================================================
%   DROPPING A QUANTUM TEMPORARY -- the paper.
%
%   Build:  make paper      full.pdf and supplement.pdf
%           make figures    figures/*.tex -> PDF (acmart never sees tikz)
%
%   Two markers say what is settled.  A blue rule and "In revision" means the
%   prose is still being reworked; a boxed "Planned" means the result belongs
%   here and is not written yet.  Skip both and you are reading the finished
%   document.  Their definitions are part 6 of preamble.tex.
%
%   Statements shared with the supplement are macros in statements.tex, so
%   the two documents cannot drift.  Appendices live in supplement.tex, a
%   separate PDF: a \cref into it does not resolve, so say "the supplement,
%   Appendix A".
%
%   Do not add a package without reading the head of preamble.tex, do not
%   change the class options, do not call \inferrule directly (use \jrule /
%   \jaxiom), and draw figures in figures/ rather than inline.
% ===========================================================================

\documentclass[acmsmall,screen,review,anonymous,nonacm]{acmart}

\usepackage{amsthm}
\usepackage{enumitem}
\usepackage{array}
\usepackage{framed}

\usepackage{mathpartir}

\usepackage{cleveref}

\setlength{\emergencystretch}{2em}
\hyphenpenalty=1000
\sloppy

\theoremstyle{acmplain}
\newtheorem{theorem}{Theorem}[section]
\newtheorem{proposition}[theorem]{Proposition}
\newtheorem{lemma}[theorem]{Lemma}
\newtheorem{corollary}[theorem]{Corollary}
\newtheorem{fact}[theorem]{Fact}

\theoremstyle{acmdefinition}
\newtheorem{definition}[theorem]{Definition}
\newtheorem{example}[theorem]{Example}
\newtheorem{remark}[theorem]{Remark}

\makeatletter
\@ifpackageloaded{cleveref}{%
  \newcommand{\crefsharedcounter}[1]{%
    \@ifundefined{#1}{}{%
      \expandafter\let\csname clv@#1\expandafter\endcsname\csname #1\endcsname
      \expandafter\let\csname endclv@#1\expandafter\endcsname\csname end#1\endcsname
      \expandafter\def\csname #1\endcsname{\crefalias{theorem}{#1}\csname clv@#1\endcsname}%
      \expandafter\def\csname end#1\endcsname{\csname endclv@#1\endcsname}%
    }%
  }%
  \crefsharedcounter{proposition}%
  \crefsharedcounter{lemma}%
  \crefsharedcounter{corollary}%
  \crefsharedcounter{fact}%
  \crefsharedcounter{definition}%
  \crefsharedcounter{example}%
  \crefsharedcounter{remark}%
}{}
\makeatother

\crefname{fact}{fact}{facts}
\Crefname{fact}{Fact}{Facts}

\renewcommand{\qedsymbol}{\ensuremath{\square}}

\definecolor{accent}{HTML}{1F4E79}
\definecolor{accentwarm}{HTML}{9A5B00}

\newcommand{\rulename}[1]{{\normalfont\scshape\footnotesize\color{accent}#1}}

\mprset{andskip=1.6em plus 0.4fil minus 0.4em, sep=0.9em}
\newcommand{\ruleskip}{\vskip 0.25\baselineskip}

\ifPDFTeX\else
  \setmonofont{Iosevka Custom}[Scale=MatchLowercase,RawFeature={-liga,-calt,-clig,-dlig}]
\fi

\newlist{steps}{enumerate}{1}
\setlist[steps]{label=$\langle$\arabic*$\rangle$,ref=$\langle$\arabic*$\rangle$,
  leftmargin=2.1em,labelsep=0.5em,align=left,itemsep=0.5pt,topsep=1pt,parsep=0pt}

\newcommand{\code}[1]{\texttt{\small #1}}

\setlength{\heavyrulewidth}{0.9pt}
\setlength{\lightrulewidth}{0.4pt}
\setlength{\aboverulesep}{3pt}
\setlength{\belowrulesep}{3pt}

\setlength{\parskip}{0pt}

\setlist[itemize]{leftmargin=1.3em, itemsep=0.25\baselineskip, topsep=0.4\baselineskip, parsep=0pt, partopsep=0pt}
\setlist[enumerate]{leftmargin=1.6em, itemsep=0.25\baselineskip, topsep=0.4\baselineskip, parsep=0pt, partopsep=0pt}
\setlist[description]{leftmargin=0pt, itemsep=0.3\baselineskip, topsep=0.4\baselineskip, partopsep=0pt, font=\normalfont\itshape}

\newcommand{\C}{\mathbb{C}}
\newcommand{\Q}{\mathbb{Q}}
\newcommand{\Z}{\mathbb{Z}}
\newcommand{\Hilb}{\mathcal{H}}
\newcommand{\ket}[1]{\lvert #1 \rangle}
\newcommand{\bra}[1]{\langle #1 \rvert}
\newcommand{\proj}[1]{\lvert #1 \rangle\!\langle #1 \rvert}
\newcommand{\id}{\mathrm{id}}

\newcommand{\den}[1]{[\![#1]\!]}

\newcommand{\emptyctx}{\cdot}

\newcommand{\drop}{\texttt{drop}}
\newcommand{\lift}{\texttt{lift}}
\newcommand{\new}{\texttt{new}}
\newcommand{\sep}{\mathrm{sep}}

\newcommand{\jty}[3]{#1 \vdash #2 : #3}

\newcommand{\jrule}[3]{\inferrule*[right={\rulename{#1}}]{#2}{#3}}
\newcommand{\jaxiom}[2]{\inferrule*[right={\rulename{#1}}]{\ }{#2}}

\newenvironment{ruledisplay}
  {\par\ruleskip\begin{mathpar}}
  {\end{mathpar}\par\ruleskip}

\definecolor{planink}{HTML}{9A5B00}
\newcounter{plancount}

\newenvironment{planned}[1]
  {\par\vskip 0.45\baselineskip%
   \stepcounter{plancount}%
   \def\FrameCommand{{\color{planink}\vrule width 1.6pt}\hspace{0.6em}}%
   \MakeFramed{\advance\hsize-\width\FrameRestore}%
   \noindent{\color{planink}\sffamily\bfseries\scriptsize Planned\hspace{0.9em}}%
   \small\ignorespaces}
  {\endMakeFramed\par\vskip 0.45\baselineskip\noindent\ignorespaces}

\newcommand{\exRunning}{%
\begin{example}[the running examples]\label{ex:three}
\leavevmode
\begin{center}\small
\begin{tabular}{@{}lll@{}}
\toprule
(a) \emph{superpose, then drop} & (b) \emph{superpose, then read} & (c) \emph{copy, measure, one joint gate} \\
\midrule
\texttt{prog a (v : [2]) : [2] \{} & \texttt{prog b (v : [2]) : [2] \{} & \texttt{prog c (v : [2]) : [2] \{} \\
\texttt{\ \ new t : [2] @ |0>;}    & \texttt{\ \ new t : [2] @ |0>;}    & \texttt{\ \ new t : [2] @ |0>;} \\
\texttt{\ \ H t;}                  & \texttt{\ \ H t;}                  & \texttt{\ \ H v;\ \ CX v t;} \\
                                   & \texttt{\ \ CX t v;}               & \texttt{\ \ p = lift[\{0\}] t;} \\
                                   &                                    & \texttt{\ \ V t v;} \\
\texttt{\ \ drop t;}               & \texttt{\ \ drop t;}               & \texttt{\ \ drop t;} \\
\texttt{\ \ return v \}}           & \texttt{\ \ return v \}}           & \texttt{\ \ return v \}} \\
\bottomrule
\end{tabular}
\end{center}
\end{example}}

\newcommand{\stmtChannel}{%
\begin{proposition}[the summed channel does not decide]\label{prop:channel}
There are two well-typed programs on a temporary qubit and a visible qubit,
each with two reachable records, that have the same summed channel $\Phi_0 +
\Phi_1$ and the same weight on every branch at every input, and opposite
verdicts on every branch.
\end{proposition}}

\newcommand{\stmtProductwise}{%
\begin{proposition}[a product at every input]\label{prop:productwise}
Let $K : \den{\Omega} \to \Hv \otimes \Hrest$ be linear, with $n = \dim\den{\Omega}$,
and suppose $Kx$ is a product vector, or zero, for every $x$. Then $K$ has one of
two shapes, and both exactly when $K$ has rank at most one.
Here the language spaces $\Hv$ and $\Hrest$ have positive dimension;
the empty register tuple denotes $\C$.
\begin{description}
\item[the temporary is fixed] $K = g \otimes W$ for a vector $g \in \Hv$
  and a linear $W : \den{\Omega} \to \Hrest$ --- the rest varies, the temporary
  does not;
\item[the rest is fixed] $Kx = Gx \otimes w$ for a fixed nonzero
  $w \in \Hrest$ and a linear $G : \den{\Omega} \to \Hv$ --- the temporary
  varies, the rest does not.
\end{description}
The first shape is what \cref{def:site} asks of a site map. A map of
the second shape is $g \otimes W$ exactly when $G$ has rank at most one, so
such an isometry with $n \ge 2$ never is.
\end{proposition}}

\newcommand{\stmtGram}{%
For a site map $K:\den\Omega\to\Hv\otimes\Hrest$, including a map on a measurement record, put
\[
 d=\dim\Hv\ge1,\qquad n=\dim\den\Omega\ge1,\qquad
 K_i=(\bra i\otimes I)K,\qquad
 \Gamma_{ij}=\operatorname{tr}(K_i^{\dagger}K_j).
\]
The Gram matrix $\Gamma$ records the pairwise inner products of the
operators $K_i$. No isometry hypothesis is needed in the proposition.
\begin{proposition}[the Gram test]\label{prop:gram}
\leavevmode
\begin{enumerate}[beginpenalty=10000]
\item \emph{Fixed block.} At a site without loans this is the safety test:
the following are equivalent:
\[
 K=g\otimes W\text{ for a unit }g
 \quad\Longleftrightarrow\quad \operatorname{rank}\Gamma\le1
 \quad\Longleftrightarrow\quad
 \Gamma^2=(\operatorname{tr}\Gamma)\Gamma.
\]
\item \emph{Least block dimension.} The map
$\rho\mapsto\sum_iK_i\rho K_i^{\dagger}$ can be represented using a
block of dimension $r\ge1$ exactly when
$r\ge\max(1,\operatorname{rank}\Gamma)$.
\item \emph{Changing gates.} A unitary on the rest leaves $\Gamma$
unchanged. A unitary $U$ on the block changes it to
$\bar U\Gamma U^{\top}$.
\end{enumerate}
\end{proposition}}

\newcommand{\stmtScalar}{%
\begin{proposition}[the block may learn nothing and the drop be unsafe]\label{prop:scalar}
Let $K=\sigma_s$ be the isometric map of a measurement-free site with
no loans in its block. Use the operators $K_i$ of \cref{prop:gram}.
The block has the same reduced state at every unit input exactly when
\[
 K_i^{\dagger}K_j=c_{ij}I\qquad\text{for all }i,j.
\]
Under this condition the drop is safe exactly when $c$ has rank one.
Then $c_{ij}=\bar g_i g_j$ for a unit certificate $g$. In contrast,
the program below has the constant mixed state $c=I/2$ and is unsafe.
\end{proposition}}

\newcommand{\stmtCliffordT}{%
\begin{theorem}[the Clifford$+T$ coefficient ring splits]\label{thm:cliffordT}
Let $\mathcal R = \Z[1/\sqrt2, i]$ and $d \ge 1$.
\begin{enumerate}
\item Every rank-one projection on $\C^{d}$ with entries in $\mathcal R$ is
  $zz^{\dagger}$ for a vector $z$ over $\mathcal R$.
\item For unit $g$, the ray $[g]$ carries a unit vector over $\mathcal R$ exactly
  when its projection $gg^{\dagger}$ has entries in $\mathcal R$.

  At a safe temporary site over $\mathcal R$, suppose the surviving weight
  at some basis input is a unit of the ring. The certificate then carries
  a unit vector over $\mathcal R$.

  A measurement-free temporary site over $\mathcal R$ has a witness over
  $\mathcal R$ exactly when it is safe: its record has weight $1$.
\end{enumerate}
\end{theorem}}

\newcommand{\stmtRings}{%
\begin{proposition}[the restricted rings]\label{prop:rings}
Over $\Z[1/2, i]$ and over $\Z[i/\sqrt2]$, every rank-one projection with
entries in the ring, in every dimension, is $zz^{\dagger}$ for a $z$ over the
ring. Over $\Z[1/2]$ the projection onto $\ket+$ is not, and over
$\Z[1/\sqrt2]$ the projection onto $\cos\tfrac{\pi}{8}\ket0 +
\sin\tfrac{\pi}{8}\ket1$ is not.
\end{proposition}}

\newcommand{\seqop}{\mathbin{+\mkern-8mu+}}
\newcommand{\Hhid}{\Hilb_{\mathsf{hid}}}

\newcommand{\stmtTypeSafety}{%
\begin{theorem}[type safety]\label{thm:typesafety}
Let $\vdash \langle \Delta, \sigma \parallel M\rangle : A$ be a well-typed
configuration of the measurement-free machine.
\begin{enumerate}
\item \emph{Progress.} Either $M$ is \texttt{return }$v$, and the
  configuration is final, or exactly one rule applies to it, and the successor is
  determined once the name \textsc{Step-New} allocates is fixed.
\item \emph{Preservation.} The step yields a well-typed configuration at
  the same type $A$.
\item \emph{Termination.} Every run from the configuration reaches a final
  configuration in exactly $|M|$ steps, where $|M|$ counts its statements
  and does not count \texttt{return}.
\item \emph{Adequacy.} The run from $\langle \Omega, x \parallel M\rangle$,
  the input registers allocated and the names it allocates fresh for $M$, ends in the
  state $\den{M}\,x$.
\end{enumerate}
\end{theorem}}

\newcommand{\stmtDesugar}{%
\begin{lemma}[desugaring]\label{lem:desugar}
Suppose $\Omega_1\vdash M_1:A$,
$\Omega_{\bar y}\vdash\bar y:A$, and
$\Omega_2,\Omega_{\bar y}\vdash M_2:B$.
The three contexts are pairwise disjoint. Every name bound in $M_1$
(register or parameter) is fresh for $M_2$, $\bar y$ and $\Omega_2$.
Every name bound in $M_2$ is fresh for $\Omega_1$ and $\bar y$.
\begin{enumerate}
\item $\Omega_1, \Omega_2 \vdash M_1 \seqop_{\bar y} M_2 : B$.
\item Every run of $M_1$ to \texttt{return }$v$, followed by the run of
  $M_2[v/\bar y]$ from the state it reaches, is a run of $M_1 \seqop_{\bar y} M_2$:
  same record and state. In the measurement-free case,
  \[
  \den{M_1\seqop_{\bar y}M_2}
  =(\id_{\Hhid^{M_1}}\otimes\den{M_2})
    \circ(\den{M_1}\otimes\id_{\den{\Omega_2}}).
  \]
\item $|M_1 \seqop_{\bar y} M_2| = |M_1| + |M_2|$.
Here $|M|$ counts statements along the longest branch of $M$:
\texttt{return} counts zero, each other statement counts one, and an
\texttt{if} counts one plus the larger count of its two arms.
\end{enumerate}
The disjointness of $\Omega_{\bar y}$ from $\Omega_1$ is not decoration:
without it the componentwise renaming can clone a register. The converse of (1) is
false --- the continuation may name a register of $\Omega_1$ directly.
\end{lemma}}

\newcommand{\stmtCompile}{%
\begin{theorem}[compilation]\label{thm:compile}
Let $M$ be a measurement-free program in the compilation fragment.
Assume each safe site $s_1,\ldots,s_r$ is emittable. Write $E_M$ for the
emitted circuit's map and $R^{s_j}$ for the ideal reclaiming gate at $s_j$.
\begin{enumerate}
\item \emph{Full state.} There is one scalar $z$, $|z|=1$, such that
\[
 E_M=z\,(R^{s_1}\otimes\cdots\otimes R^{s_r}\otimes I)\circ\den M.
\]
The output factors are grouped by dropped block, in site order, followed
by the remaining factors. The scalar $z$ is independent of every input,
including the states of the loans; it is the product of the phases
omitted during emission.
\item \emph{Returned output.} The emitted circuit and $M$ have the same
reduced state on the returned registers at every input. If no reclaiming
gate is inserted, their full maps differ only by $z$.
\item \emph{Reusable registers.} Each reclaimed temporary is in its
allocation state. Each reclaimed loan is in its original state, with
its correlations with the caller preserved. These are equalities of
states, so the global phase in (1) does not affect them.
\end{enumerate}
\end{theorem}}

\newcommand{\Hv}{\Hilb_{v}}
\newcommand{\Hrest}{\Hilb_{\mathsf{rest}}}
\newcommand{\Hloan}{\Hilb_{\mathrm{loan}}}
\newcommand{\Hother}{\Hilb_{\mathsf{other}}}
\newcommand{\pfx}[2]{#1|_{#2}}

\newcommand{\stmtSite}{%
\begin{definition}[the site, and safety at it]\label{def:site}
Let $\jty{\Omega}{M}{A}$ be a program.
\begin{enumerate}[beginpenalty=10000]
\item A \emph{site} $s$ of $M$ is an occurrence of $\drop\,v$ in it. Its
  \emph{block} is the ordered list of registers $v$ names --- none for $v = ()$, a site
  with $\Hv = \C$, safe with the unit scalar as certificate;
  $\Hv$ is the tensor product of
  their spaces in that order. The space $\Hrest$ is the tensor product of all remaining
  register spaces: inputs in context order, followed by allocations,
  including previously dropped registers (\cref{sec:lang:dynamics}).
\item The \emph{site map} $\sigma_s : \den{\Omega} \to \Hv \otimes \Hrest$
  sends an input vector $\xi\in\den\Omega$ to the joint state of all
  registers just before the drop. The block factors are listed first;
  no physical gate is applied by this regrouping.
\item First suppose none of the registers named by this drop is marked
  \texttt{loan}. They may be allocated temporaries or unmarked inputs.
  The site is \emph{safe} when its site map
  factors,
  \[ \sigma_s = g \otimes W, \]
  for a unit vector $g \in \Hv$, the site's \emph{certificate}, and a linear
  $W : \den{\Omega} \to \Hrest$: the block comes out in the one state $g$
  whatever the input.
\item If the drop also names registers marked \texttt{loan}, they must be
  returned in their unknown original states. Group those input factors as $\den{\Omega} =
  \Hloan \otimes \den{\Omega'}$, the block's loans and the rest, and the site
  is \emph{safe} when
  \[ \sigma_s = B \otimes W \]
  for an isometry $B : \Hloan \to \Hv$, the site's \emph{block factor}, and a
  linear $W : \den{\Omega'} \to \Hrest$. Clause (3) is the case $\Hloan = \C$,
  $B = g$.
\end{enumerate}
\end{definition}}

\newcommand{\stmtModular}{%
\begin{proposition}[safety is modular]\label{prop:modular}
Let $\sigma_s$ be the linear map of a site as in \cref{def:site}.
Let $\Hilb_c\ne0$ be a client register on which the program does nothing.
Then
\[
 \sigma_s\otimes I_c=B\otimes W'
 \text{ for some }W'
 \quad\Longleftrightarrow\quad
 \sigma_s=B\otimes W\text{ for some }W.
\]
The same isometric block factor $B$ works on both sides, including when
the client is entangled with the inputs. For a block without loans,
read $B=g$.
\end{proposition}}

\newcommand{\stmtIsoExt}{%
\begin{lemma}[isometry extension]\label{lem:isoext}
Let $\Hilb$ be a finite-dimensional complex Hilbert space. Every isometry
$V : S \to \Hilb$ out of a subspace $S \subseteq \Hilb$ is the restriction of
a unitary of $\Hilb$.
\end{lemma}}

\newcommand{\stmtCriterion}{%
\begin{theorem}[the criterion, and contextual equivalence]\label{thm:criterion}
In the Setting, the following are equivalent.
\begin{enumerate}
\item $s$ is safe.
\item $\sigma_s = g \otimes W$ for a unit $g \in \Hv$ and an isometry
  $W : \den{\Omega} \to \Hrest$.
\item Some weakening witness $\sep(R,W)$ has denotation $\sigma_s$.
\item No closed context tells $\pfx{M}{s}$ from some $\sep^{\circ}(R,W)$.
\end{enumerate}
Two prefixes no closed context tells apart have the same certificates.
\end{theorem}}

\newcommand{\stmtComposite}{%
\begin{proposition}[the composite temporary]\label{prop:composite}
Let $M$ have temporaries $t_1,\ldots,t_n$, each dropped once on the
chosen record, and let $\Phi$ be its final state map on that record.
For clause (1), assume that the tail after the drop of $t_k$ is
injective on the span of the row vectors
$(\bra{e_i}\otimes I)\sigma_s(x)$, as $i$ and $x$ vary. This holds if
that tail contains no lifts. For a lift, the condition says that its
projection kills no nonzero vector in the relevant span.
Call $t_k$ \emph{separately
constant} when $\Phi(x)$ lies in $g_k \otimes (\text{the other factors})$ for
one unit $g_k \in \Hilb_{t_k}$ and every input $x$.
\begin{enumerate}
\item The site of $\drop\,t_k$ is safe if and only if $t_k$ is separately
  constant.
\item Every $t_k$ is separately constant if and only if the composite site
  $\drop\,(t_1, \ldots, t_n)$, placed at the end after deleting the
  individual drops, is safe with
  a certificate carrying the product $g_1 \otimes \cdots \otimes g_n$.
\end{enumerate}
\end{proposition}}

\newcommand{\stmtPhase}{%
\begin{lemma}[equal answers, proportional maps]\label{lem:phase}
Let two prefixes with every $\drop$ deleted have the same context and type,
so that their denotations $\sigma, \sigma' : \den{\Omega} \to \Hilb$ are
linear maps into one space. If no closed context tells them apart, then
$\sigma' = c\,\sigma$ for a scalar $c$, of modulus one once $\sigma \neq 0$.
Conversely, no closed context tells apart two prefixes with
$\sigma' = c\,\sigma$ and $|c| = 1$.
\end{lemma}}

\newcommand{\stmtReuse}{%
\begin{proposition}[reclaiming temporaries]\label{prop:reuse}
In the Setting above, define
\[
 R^{\psi}=\begin{cases}
 I,&w=0,\\
 I-2ww^{\dagger}/\lVert w\rVert^2,&w\ne0.
 \end{cases}
\]
\begin{enumerate}
\item \emph{Exchange.} $R^{\psi}$ is the unique linear map that exchanges
$\tilde g$ and $\psi$ and fixes their common orthogonal complement.
It is a unitary involution.
\item \emph{No gate needed.} $R^{\psi}=I$ exactly when $\tilde g=\psi$,
equivalently when $g$ lies on the allocation ray.
\item \emph{Restoration.} At every input, applying $R^{\psi}$ to the
block puts it in $\psi$ and leaves the reduced state of the rest unchanged.
\end{enumerate}
\end{proposition}
A dropped block may contain both loans and temporaries. A loan must be
returned in the state the caller supplied; a temporary must be returned
to its declared allocation state.

At a safe site, the block factor
$B:\Hloan\to\Hv$ of \cref{def:site}(4) describes where the program has
put the loans' information. For a loan input $\ell$, the block contains
$B\ell$. Let $J$ describe the desired result: $J\ell$ places the unchanged
loans beside the temporaries in their allocation state $\psi$.
If loans precede temporaries in the tuple, $J\ell=\ell\otimes\psi$;
otherwise these factors are placed in the tuple's order.

\begin{proposition}[returning loans]\label{prop:loanreturn}
For the safe block above, there is a unitary $T$ on the block such that
\[
                    T(B\ell)=J\ell\qquad\text{for every loan input }\ell.
\]
Thus the same gate restores every possible loan state and every
temporary's allocation state. It also preserves the loans' correlations
with registers kept by the caller.

Only the action of $T$ on the image of $B$ is prescribed. Different
unitaries may extend that action to the rest of the block space.
If the block contains only loans, $B$ is square and unitary:
$J=I$ and $T=B^{\dagger}$. In this case $T$ is the identity up to phase
exactly when $B$ is.
\end{proposition}}

\newcommand{\stmtReuseSound}{%
\begin{proposition}[reuse is sound]\label{prop:reusesound}
Extend the machine of \cref{fig:machine} with a pool: at a safe site,
\textsc{Step-Drop-Reclaim} applies $R^{s}$ to the block and moves its
registers to the pool instead of hiding them (loans are returned by the gate of \cref{prop:loanreturn}
and are never pooled);
\textsc{Step-New-Reuse} serves
a \texttt{new } $t : [d] \mathbin{@} \psi$ from a pooled register standing in $\psi$
instead of allocating a fresh one.
\begin{enumerate}
\item Every run of the pooled machine is simulated by a run of the machine of
  \cref{fig:machine}: the same records, the same weight at every input, and
  the same reduced state on the returned registers, on every record.
\item The pooled machine allocates at most as many registers, and strictly
  fewer whenever a reclaimed register is served.
\item If reclamation is allowed at an unsafe site, reuse can be unsound. Extend program (b) by
  allocating a fresh qubit and returning it with $v$. If that allocation
  instead reuses the dropped block, no gate on the block alone makes the
  returned joint state agree at every input.
\end{enumerate}
\end{proposition}}

\newcommand{\stmtTypeSafetyLift}{%
\begin{theorem}[type safety, with the two forms]\label{thm:typesafetylift}
\Cref{thm:typesafety} holds of the five-rule machine, with two clauses
read again and two unchanged.
\begin{enumerate}[beginpenalty=10000]
\item \emph{Progress} now reads: a well-typed configuration is final, or
  exactly one rule applies to it, and at \textsc{Step-Lift} both of that rule's
  labelled instances do --- a lift being the only labelled step.
\item \emph{Termination} now reads along the branch taken: the run whose
  record is $o$ reaches a final configuration in at most as many steps as that branch
  has statements, an \texttt{if} counting as one, and so in at most $|M|$.
\item \emph{Preservation} stands word for word, and so does the
  hidden-register invariant of \cref{sec:lang:dynamics}: a lift binds its parameter
  and leaves its register live under its own name, so no register enters or leaves
  the context there.
\item \emph{Adequacy} stands with the run read along its record: the run
  from $\langle \Omega, x \parallel M\rangle$ whose record is $o$ ends in the state
  $E_o(M)\,x$.
\end{enumerate}
\end{theorem}}

\newcommand{\stmtBlockLift}{%
\begin{remark}[a block lift is a derived form]\label{rem:blocklift}
Let $S\subseteq[d]$ be nonempty and proper. Define the indicator gate
\[
 U_S=\sum_{j\in S}\proj j\otimes X+
     \sum_{j\notin S}\proj j\otimes I.
\]
It leaves $x$'s level unchanged and flips a fresh qubit $a$ exactly when
that level belongs to $S$. In a signature containing $U_S$, replace
$p=\lift_S\,x\texttt{;}\ M$ by these steps:
\begin{enumerate}
\item Allocate $a$ in $\ket0$ and apply $U_S$ to $(x,a)$.
\item Set $p=\lift_{\{1\}}\,a$. Thus $p=1$ has the same meaning as in
      the original test: the level of $x$ belongs to $S$.
\item Drop $a$ and continue with $M$.
\end{enumerate}
To check the replacement, write the incoming state as
$\sigma=\sum_j\ket j_x\otimes\sigma_j$ and put
$\tau_1=\sum_{j\in S}\ket j\otimes\sigma_j$ and
$\tau_0=\sum_{j\notin S}\ket j\otimes\sigma_j$.
Step (1) gives $\tau_1\otimes\ket1_a+\tau_0\otimes\ket0_a$.
On outcome $b$, step (2) keeps
$\tau_b\otimes\ket b_a=((P_S^b)_x\sigma)\otimes\ket b_a$.
This is the original branch state with one extra, fixed factor.
The replacement therefore has the same type and per-record map after
omitting that factor. Its drop of $a$ is safe, and $X^b$ restores $a$
to $\ket0$ (\cref{prop:reuse}).
\end{remark}}

\newcommand{\stmtReuseRecord}{%
\begin{proposition}[the record chooses the reclaiming gate]\label{prop:reuserecord}
Let $s$ fire on record $o$ with $\sigma_s^o\ne0$, safe there with certificate
$[g_o]$ and block allocated in $\psi$, and let $R_o$ be the reclaiming gate
of \cref{prop:reuse} for that certificate.
\begin{enumerate}
\item $g_o$ and $R_o$ depend on $o$ only through the bits $o$ takes at the
  lifts that fire before $s$: two records agreeing on those have the same
  certificate and the same reclaiming gate.
\item One gate cannot serve all records: some well-typed $M$ has a
  site whose reclaiming gate is $I$ on one record and $X$ on the other,
  whether or not an \texttt{if} of $M$ reads the bit that decides.
\end{enumerate}
\end{proposition}}

\newcommand{\stmtCompileLeaf}{%
\begin{theorem}[compilation, one record at a time]\label{thm:compileleaf}
Let $M$ be in the compilation fragment, with every lifted register of
power-of-two dimension. Fix a record $o$. List the safe sites that fire
on this record as $s_1,\ldots,s_r$, in execution order. Assume that each
site's certificate or loan block factor is emittable, and write
$R_o^{s_j}$ for its ideal reclaiming gate.

Let $C_o$ be the corresponding branch map of the emitted circuit.
When comparing maps, omit any additional indicator qubits that the
emitter returns to $\ket0$. Order the remaining hidden factors by site
and drop-tuple order, as in \cref{thm:compile}. Then:
\begin{enumerate}
\item The compiled branch is the source branch followed by the chosen
  reclaiming gates:
  \[
  C_o=z_o\bigl(R_o^{s_1}\otimes\cdots\otimes R_o^{s_r}\otimes\id\bigr)E_o(M),
  \qquad |z_o|=1.
  \]
  The phase $z_o$ is independent of every input, including loan states.
\item The branch weight and the reduced output state agree with those
  of $E_o(M)$ for every input. If no gate is emitted, the entire
  branch maps agree up to $z_o$.
\item Each reclaimed temporary is in its declared allocation state,
  and each reclaimed loan is in its original state. The phase in
  clause (1) changes neither of these states.
\end{enumerate}
\end{theorem}}

\newcommand{\stmtRecordWitness}{%
\begin{theorem}[the criterion, one record at a time]\label{thm:perrecord}
Let $s$ be a site of $M$ whose block is a tuple of temporaries allocated in
the unit state $\psi$, and $o$ a record on which $s$ fires. The following are equivalent.
\begin{enumerate}
\item $s$ is safe on $o$.
\item $\sigma^{o}_s = g \otimes W_o$ for a unit $g \in \Hv$ and a contraction
  $W_o : \den{\Omega} \to \Hrest$.
\item A witness has map $\sigma_s^o$ on its \emph{kept record},
where its fresh qubit $a$ reads $0$, after omitting that fixed
$\ket0_a$ factor. The witness performs three steps:
prepare the block in $\psi$ and apply one gate $R$ to it;
allocate the other temporaries and $a$ in $\ket{0}$ and apply one gate $U$
to everything except the block; then lift $a$ against $\{0\}$.
\item No closed context, and already no lift-free one, tells $\pfx{M}{s}$ on
  $o$ from some such program on its kept record.
\end{enumerate}
\end{theorem}}

\newcommand{\stmtObsInvRecord}{%
\begin{corollary}[safety is an observational invariant, one record at a time]
\label{cor:obsinvrecord}
Fix a record $o$ and let two prefixes with every $\drop$ deleted have the
same context and type, and let no lift-free closed context tell them apart on $o$. Then
${\sigma'}^{o} = c\,\sigma^{o}$ for a scalar $c$, of modulus one once
$\sigma^{o} \ne 0$; hence their sites are both safe or both unsafe on $o$, and every
certificate of one is a certificate of the other.
\end{corollary}}

\newcommand{\stmtNoTextQubit}{%
\begin{theorem}[the verdict is record-dependent]\label{thm:notext}
Let $c$ be program (c) of \cref{sec:meas:programs}.
\begin{enumerate}
\item $c$ is well typed, and its two records are both reachable, each of
  weight $\tfrac12$ at every basis input.
\item Under the record taking $1$ at the lift, the temporary landed
  in $\{0\}$, and the site map is
  $\sigma_{\mathrm{in}}(x) = \tfrac{x_0 + x_1}{\sqrt2}\,\ket{0} \otimes \ket{0}$,
  and the drop is safe, with certificate $\ket{0}$.
\item Under the record taking $0$ the site map is
  $\sigma_{\mathrm{out}}(x) = \tfrac{x_0 - x_1}{2}\,(\ket{0} \otimes \ket{1}
   + \ket{1} \otimes \ket{0})$,
  and the drop is unsafe.
\end{enumerate}
\end{theorem}}

\newcommand{\stmtRelativisedSite}{%
\begin{theorem}[the criterion over a gate ring]\label{thm:relativised}
Let $\mathcal A$ be a gate ring, $M$ a measurement-free program over
$\mathcal A$, and $s$ a site of $M$ whose block is a tuple of temporaries
allocated in $\psi \in \mathcal A^{d}$.
\begin{enumerate}
\item If some weakening witness $\sep(R,W)$ with $R$ over $\mathcal A$ has
  denotation $\sigma_s$, then $s$ is safe with a certificate carrying the unit
  vector $R\psi$ over $\mathcal A$, whatever $W$ is.
\item If $\mathcal A$ has a column lemma (\cref{fact:column}) and $s$ is safe
  with a certificate carrying a unit vector over $\mathcal A$, then some
  weakening witness \emph{over $\mathcal A$} has denotation $\sigma_s$.
\end{enumerate}
\end{theorem}}

\newcommand{\stmtClassicalCostSite}{%
\begin{theorem}[the classical price]\label{thm:classicalcost}
Deciding whether a given site whose block contains no loans is safe on a
given record is coNP-complete, for programs whose gates are drawn from $\{X, \mathrm{CNOT},
\mathrm{Toffoli}\}$ on qubit registers and whose allocations are clean.
Hardness holds already for measurement-free programs with one site.
\end{theorem}}


\settopmatter{printacmref=false, printfolios=true}
\setcopyright{none}
\acmConference[Under review]{Submitted for double-blind review}{}{}
\acmYear{2026}

\begin{document}

\title{Dropping a Quantum Temporary}

\author{Author Name}
\orcid{0000-0000-0000-0000}
\affiliation{%
  \institution{Indiana University}
  \city{Bloomington}
  \state{Indiana}
  \country{USA}}
\email{sdkadu@iu.edu}

\begin{CCSXML}
<ccs2012>
<concept>
<concept_id>10011007.10011006.10011039</concept_id>
<concept_desc>Software and its engineering~Formal language definitions</concept_desc>
<concept_significance>500</concept_significance>
</concept>
<concept>
<concept_id>10003752.10003790.10003806</concept_id>
<concept_desc>Theory of computation~Programming logic</concept_desc>
<concept_significance>300</concept_significance>
</concept>
</ccs2012>
\end{CCSXML}
\ccsdesc[500]{Software and its engineering~Formal language definitions}
\ccsdesc[300]{Theory of computation~Programming logic}

\keywords{quantum programming languages, uncomputation, ancilla, deallocation,
dynamic lifting, type systems, contextual equivalence}

\maketitle


% Section 1 -- Introduction
\section{Introduction}\label{sec:intro}

Here is a simple two-gate quantum circuit. It allocates a qubit, applies a Hadamard,
applies a second Hadamard, and asks the compiler to forget it:

\begin{center}
\includegraphics{figures/two-gate}
\end{center}

\emph{Forget} here means something specific, and it is the subject of
the paper: we want the qubit to go back to a free pool, and a later
allocation should be able to \emph{reuse} it without reinitialising
it. Qubits are the scarce resource. A compiler that cannot reuse a
temporary allocates a fresh one, and the width of the program grows
with every scratch value it ever computed. The established route to
reuse is \emph{uncomputation}: run the computation that produced the
scratch value backwards, so that the qubit is provably back at
$|0\rangle$, and hand it to the pool clean. That is one way to make
reuse safe. This paper asks what reuse actually requires, and the
answer is, perhaps surprisingly, weaker than restoration: the qubit
need not be back at $|0\rangle$, but it must carry no information the
rest of the program can read. What we certify is the state in which a
temporary may be released; the allocator itself, and what it means to hand the
same qubit out again under a measurement record, are outside the model.

Every classical model of computation throws things away for free---a
Turing machine overwrites its tape, a $\beta$-reduction discards its redex,
a stack pops---because overwriting is a primitive irreversible operation
there, and the classical alternatives it hides cannot interfere. Bennett's
reversible machine buys the discarded value back, ending with its output
beside ``a reconstructed copy of the input''~\cite{Bennett1973Logical}. A
superposition of control withdraws the licence. The temporary is a tensor
factor of the state space and not one alternative among several, so discarding
it is not the choosing of a branch: the joint state pairs each temporary basis
vector with a visible vector of its own, and the inner products among those
decide what interference survives. That is why \emph{forget} has content here, and why its
condition cannot be read off how the value was produced.

In Silq~\cite{Bichsel2020Silq} the circuit is the program:
\begin{center}
\begin{tabular}{@{}p{0.34\linewidth}@{\qquad}p{0.52\linewidth}@{}}
\begin{minipage}[t]{\linewidth}
\centering\ttfamily
\begin{tabular}[t]{@{}l@{}}
t0 := 0:B;\\
t1 := H(t0);\\
t := H(t1);\\
forget(t);
\end{tabular}
\end{minipage}
  &
    \small The register is $|0\rangle$: two Hadamards compose to the identity.
    Discarding it is safe, and Silq agrees that it is safe if we use the variant
    \code{forget(t=0:B)}. The automatic \code{forget}, however,
    is refused with error: \emph{cannot synthesize forget expression.}
\end{tabular}
\end{center}

The refusal is not really a bug: Silq's rule asks how the value
was \emph{produced}, not what the register \emph{is}. For Silq, $H$ is
not \code{qfree}, which is Silq's word for an operation whose
semantics ``can be described classically'' and which never creates
superposition~\cite[\S4.1]{Bichsel2020Silq}, so no uncomputation can
be synthesised for \code{t}, and the automatic \code{forget} is
refused. Unqomp refuses it as well for the same reason: it synthesises
uncomputation rather than checking it, and its correctness section
shows why a non-\code{qfree} gate on an ancilla defeats the
mechanism~\cite{Paradis2021Unqomp}. The other kind of existing check
accepts this program, for an instructive reason. Q\# requires a
released qubit to be back in the state it was allocated
in~\cite{Svore2018Qsharp,MicrosoftQsharpMemory}, and Quipper's
\code{with\_ancilla} closes on the programmer's promise of the
same~\cite{Green2013Quipper}. Here the register \emph{is} back at
$|0\rangle$. Program (a) below is where those rules refuse a safe
drop. Every rule of either kind draws a line through the safe
programs, and the systems that publish their rule say so: ReQWIRE
calls its syntactic conditions ``sufficient but not
necessary''~\cite{Rand2019ReQWIRE}, and Qutes calls its own detection
``necessarily incomplete''~\cite{Faro2026Qutes}. Silq is the one we
ran, and \cref{sec:related} compares the rest.

So the gap is known and in print. \emph{The gap:} a drop is safe when
no continuation of the program can observe it, that is, when the
temporary carries no information about the visible register at the
moment it is dropped. The existing checks ask about the program's
syntax instead---how the value was produced, whether the register is
back where it started, who owns it and for how long---and
\cref{sec:overview} sets the three against each other. Neither of the
first two is complete, and the two programs above are the witnesses:
the two-gate circuit is a safe drop the producer check refuses, and
program (a) is a safe drop the restoration check refuses. The recent
tools decide a neighbouring predicate at circuit level: whether a
circuit restores its ancilla~\cite{Li2026AncillaSafety}, whether a
restoring uncomputation exists~\cite{Liu2026Uncomputation}, and how to
schedule borrows given blocks already safe~\cite{Xu2026Bona}. Drop
safety as we formulate it here lies strictly between restoration and
restorability. Call a program \emph{restored} when it itself returns the
temporary to the state it was allocated in, and \emph{restorable} when some
continuation could~\cite{Liu2026Uncomputation}. In the unrestricted
measurement-free calculus
\emph{restored} implies \emph{safe to drop} implies \emph{restorable},
and neither converse holds: program (a) is safe and not restored, and
Bennett's compute--copy--uncompute, with its target counted into the
temporary, is restorable and, unless its circuit is constant, not safe
(\cref{sec:cost}). Over a fixed coefficient ring the second
implication asks for more---the splitting condition of
\cref{sec:gateset}, and over a finite repertoire an exact-synthesis
step besides. None of those three tools reads a program with
measurement, where the verdict may depend on the measurement record
(\cref{sec:boundary}). We ask two questions about the gap. \emph{What
  is the condition, stated as a property of programs rather than of
  matrices?} And \emph{can the gap be closed at all?} The first has an
exact answer. The second has an exact answer too, and it is
negative---but only past a boundary we can draw, and the negative has
a type that the natural way of stating it gets wrong.

\subsection*{Three programs}

We work in a small language with one scoped temporary $t$, one visible
register $v$, and an affine $\drop$ whose typing rule carries no side
condition. A side condition would
have to say \emph{this temporary is uncorrelated from $v$}, a fact
about what the program computes, and a rule that demanded it would
make the unsafe programs untypable, so the checks that refuse them
could not be compared.

The \emph{temporary} $t$ is the ancilla of the quantum computing
literature: a register allocated inside a scope and dropped at its
end. The \emph{visible} register $v$ is everything else in scope---the
data qubits---and it is called visible because it is what the rest of
the program can see after the drop. Both are \emph{registers} of
arbitrary finite dimension, so ``one temporary and one visible
register'' is not two qubits unless the dimension is fixed at two. An $n$-qubit
workspace is a visible register of dimension $2^n$, and several
temporaries in nested scopes compose into one big temporary of product
dimension with the verdict preserved exactly (\cref{sec:many}). The
results therefore apply to programs of any width, and \cref{sec:eval}
runs the checker on eight textbook algorithms. \Cref{sec:gateset} is the one
section needing a ring-theoretic input, the column lemmas of the
exact-synthesis literature, and it says which rings have them. One
temporary is the smallest system in which the question has content:
without the visible register there is no observer, and without joint
gates there is no channel. An impossibility is only as strong as its
hypotheses are weak, so the negatives are proved there on purpose.

Each program below allocates the temporary, computes, and drops it.
A run leaves the temporary \emph{constant} when it leaves it in one fixed
state independent of the input.

\exRunning

Program (a) superposes the temporary and drops it. Its final state is
$|{+}\rangle$ whatever the input: constant, uncorrelated with $v$, and not
$|0\rangle$.

Program (b) reads the temporary into $v$ before dropping it. It
is the one-line proof that the affine rule cannot be silent: after the read,
$v$ and $t$ are entangled, and forgetting $t$ decoheres $v$.

A gate in this language is any orthogonal matrix
  (\cref{def:calculus}), and the $\mathit{tilt}$ gate in program (c) is
  $\mathit{tilt} = \tfrac13\bigl(\begin{smallmatrix} 2 & 2 & 1 \\ 2 & -1 & -2 \\ 1 & -2 & 2
\end{smallmatrix}\bigr)$, with $\mathit{tilt}\,|0\rangle = \tfrac13(2,2,1)$.

Program (c) measures the temporary, then reads it into $v$ and drops
it, and no gate after the measurement branches on the outcome. The
gate $\mathit{tilt}$ is an orthogonal gate on a three-dimensional
temporary, chosen so that the measured block $S = \{0,1\}$ and its
complement both survive with nonzero weight. Which block survives is
decided at run time. On one outcome the drop is safe, on the other it
is not, and nothing in the program's syntax says which run is
happening.

\subsection*{Why discard is not free here}

In some languages this question has no content. In the quantum $\lambda$-calculus quantum data is affine
and discard is a typing rule that asks
nothing~\cite{SelingerValiron2006}, because discard is a partial trace
over a density-matrix semantics: the continuation receives the reduced
state of the registers it can still name, which is correct whatever
the discarded register held, and that register is never reused. Here
we want it reused. The compiler traces nothing, and it hands
the same physical qubit to the next allocation, on the promise that
the next allocation will find nothing in it. \Cref{sec:calculus} makes
the promise precise: \textsc{Drop} has no side condition, the
elaboration of $\drop$ is the identity, and the result type of a
program, which no longer mentions the temporary, asserts that no
continuation will look at it. So the promise is the affine rule's
\emph{premise}, not its conclusion, and the rule leaves it unchecked
on purpose: the programs that break it must be expressible, as
program (b) of \cref{ex:three} is, well typed with a wrong drop.

\subsection*{What a compiler gets, and what it pays}
The criterion is \emph{stated} on programs and \emph{decided} on
matrices, and the two are not in tension. Stated: a drop is safe when
the program is equivalent to one that never touched the temporary, a
sentence with no matrix in it (\cref{sec:criterion}). Decided: the
program text names its gates, each gate is a matrix, and multiplying
them in the order the text gives yields the denotation $K$
(\cref{sec:calculus}). The decision then is a rank test on $K$,
static, needing neither an input nor a run. The price is a
proportionality test on the $d$ Kraus operators, each an $m \times m$
matrix: polynomial in the joint dimension $dm$, hence exponential in
the qubits in scope, the cost of writing the denotation down at all.
Once the verdict depends on which outcome occurred, one verdict for
all records is wrong on one of them, and a sound one is wrong by refusing a
safe run (\cref{sec:boundary}). From
\cref{sec:residual} a compiler gets instead a table with one row per
measurement record, at most $2^{\ell}$ for $\ell$ lifted measurements
on the command's longest path, each row one proportionality test,
computed before the run and consulted at run time by a lookup on the
classical outcomes, and $\ell$ is at most two in the programs of
\cref{sec:eval}.

\subsection*{Contributions}

\begin{itemize}
  \item \textbf{Measurement-free: an exact criterion on the program text}
        (\cref{sec:criterion}). Three conditions on a typed program coincide:
        it is \emph{weakenable}, some typed program that prepares the temporary
        once and never touches it again having the same denotation; it is
        \emph{drop-safe}, its temporary ending in one fixed state whatever the
        input; and no continuation can distinguish it from a program of that
        shape (\cref{thm:criterion,thm:ctxequiv}). All three say the channel
        the drop induces has Kraus rank one (\cref{thm:bridge}), a finite
        computation on the matrix the program's gates multiply out to.
        The rank-one arithmetic is textbook---the Choi rank of an \emph{isometric channel}~\cite{Choi1975CP,Kraus1983States,Watrous2018TQI}---and what is ours is that it is asked of a program's denotation.
        Restoring the temporary to where it was allocated is this criterion
        pinned to one ray, and it refuses eight of the twelve safe catalogue
        programs, Deutsch--Jozsa among them.
  \item \textbf{With measurement: a verdict fixed before the run is
        incomplete} (\cref{sec:boundary}). We extend the calculus with
        measurement, following Proto-Quipper's \emph{dynamic
        lifting}~\cite{Fu2023DynamicLifting}; a typed program of the affine
        calculus elaborates to a lift-free command of it, with the same
        denotation. Within the extension we exhibit a well-typed command
        with two reachable measurement records of opposite verdict on the
        same drop, and no \emph{sound} record-oblivious verdict keeps the
        safe one: the commands on which a safe run must be lost are exactly
        the outcome-dependent ones (\cref{sec:boundary}), which the presence
        of a lift does not characterise.
  \item \textbf{One exact repair: a record-indexed table}
        (\cref{sec:residual}). The boundary forces record awareness---the
        missing bit exists only after the measurement---and one exact
        realisation is a table emitted before the run, one row per record and
        at most $2^{\ell}$ rows for $\ell$ lifts, each row the proportionality
        test of \cref{sec:criterion}: the static analysis's own cost, once per
        record. The table is exact on the very command where a sound
        record-oblivious verdict must lose a safe run.
\end{itemize}

\begin{outlinebox}{gate arithmetic, and the price of the verdict}
\Cref{sec:gateset,sec:cost}, in summary. Over a gate ring completable in the
temporary's dimension the criterion acquires one side condition, that the
projector onto the final state factor as $zz^{\dagger}$ over the ring: automatic
over the Clifford$+T$ ring, false over its real subring. Deciding is
coNP-complete already on permutation programs.
\end{outlinebox}

\begin{outlinebox}{composition, lifetimes, and the implementation}
\Cref{sec:many,sec:eval}, in summary. Composition reduces several temporaries
to one, though dropping them one at a time can cost a square root; typing
splits per lifetime and drop safety does not. Two decision procedures sharing
no code but one arithmetic backend agree on every catalogue verdict.
\end{outlinebox}


% Section 2 -- The idea, on three programs
\section{The idea, on three programs}\label{sec:overview}

This section carries the three programs of \cref{ex:three} through the
whole paper informally.

\subsection{The temporary is the scrutinee of an invisible switch}

From a language implementation perspective the whole idea is that a
quantum program with a temporary of dimension $d$ is a \code{switch}
with $d$ branches. The scrutinee is the temporary's basis state: the
$i$th branch is what the program does to the visible register when the
temporary is in state $|i\rangle$. Call that branch's action $K_i$. It
is an $m \times m$ matrix on the visible register, read off the
program's denotation $K$ as the block of rows labelled $i$:
\[
  K v \;=\; \sum_{i=0}^{d-1} |i\rangle \otimes K_i v ,
  \qquad K_i = (\langle i| \otimes I)\,K .
\]
The $K_i$ are the \emph{Kraus operators}\revisinginline{~\cite{Kraus1983States}} of the channel
that dropping the temporary induces on the visible register: tracing the temporary out of
$K\rho K^{\dagger}$ leaves
$\rho \mapsto \sum_i K_i \rho K_i^{\dagger}$. Quantum mechanics adds
one thing to the switch: every branch runs, and the results are kept
side by side, each labelled by its branch. The label captures the
difficulty: the temporary holds a memory of what happened to the
visible register, and a memory is exactly what a drop must not throw
away while anyone downstream can read it.

Now $\drop\ t$ deletes the scrutinee, and the question is whether
anybody downstream can tell. If every branch does the same thing,
which in matrices means $K_i = g_i\,W$, one visible action $W$ scaled
differently per branch, then the switch was fake: the program's
behaviour never depended on the scrutinee, deleting it changes
nothing, and thus the drop is safe. If two branches do
\emph{different} things, the switch was real: the downstream code
behaves one way when the temporary went left and another when it went
right, so the temporary recorded something about the visible input,
and deleting it destroys a correlation that was there. So the safety
of a drop is a proportionality condition: the Kraus operators
$K_0, \ldots, K_{d-1}$ are pairwise proportional, all scalar multiples
of one isometry $W$. Equivalently, the channel the program induces has
\emph{Kraus rank one}: it is conjugation by that one isometry. On the matrix
a program's gates multiply out to, that is a finite calculation in exact
arithmetic, so the verdict is decided and not estimated.

The arithmetic is not new. That a channel is
$\rho \mapsto \sum_i K_i \rho K_i^{\dagger}$, that two families of operators
describe one channel exactly when an isometry relates them---a unitary, once
the shorter family is padded with zeros---and that the fewest a channel needs
is the rank of its Choi matrix, the \emph{Choi rank}, are the standard
structure theory of channels~\cite{Kraus1983States,Choi1975CP,Watrous2018TQI},
and rank one is its smallest case, the \emph{isometric channel}. What is new is what the rank is asked about: the
denotation of a \emph{program}.

\begin{figure}
\centering
\includegraphics{figures/criterion}
\caption{The criterion. A triangle starts a register in a fixed state and
the bars end one, which is $\drop$.\newline
\textbf{Left:} any program of the calculus.\newline
\textbf{Middle:} the drop-safe shape (\cref{def:core}): the temporary comes
out in one state $g$ whatever $v$ was, and nothing crosses between the
registers.\newline
\textbf{Right:} the weakenable shape (\cref{def:sw}): the temporary is
prepared by one gate $R$ and dropped, never having met $v$.\newline
The left program is drop-safe exactly when it equals the middle, and
\cref{thm:criterion} says that is exactly when it equals the right.}
\Description{Three circuit diagrams joined by equals signs: an arbitrary
two-register unitary, then a product of a fixed temporary state with a
visible gate, then a temporary prepared by one gate beside a visible gate.}
\label{fig:criterion}
\end{figure}

\subsection{Program (a) is safe, and it does not restore anything}
Program (a) applies $H$ to the temporary and nothing else. Program (b)
then reads the temporary into $v$ by \textsc{cnot}. Write each
denotation as a matrix, rows indexed by the temporary's level first, so
that the rows at level $0$ form $K_0$ and those at level $1$ form $K_1$:
\begin{center}
\small
\begin{minipage}[t]{0.47\linewidth}
\raggedright
Program (a) sends $v$ to $\tfrac1{\sqrt2}(|0\rangle + |1\rangle) \otimes v$:
\par\smallskip\centerline{$
  K_{\mathrm{(a)}} \;=\; \tfrac1{\sqrt2}
  \begin{pmatrix}
    {\color{accent}1} & {\color{accent}0} \\
    {\color{accent}0} & {\color{accent}1} \\
    {\color{accentwarm}1} & {\color{accentwarm}0} \\
    {\color{accentwarm}0} & {\color{accentwarm}1}
  \end{pmatrix}
  \begin{array}{@{}l@{}}
    \left.\begin{matrix}\phantom{0}\\ \phantom{0}\end{matrix}\right\}
      {\color{accent}K_0 = \tfrac1{\sqrt2} I} \\
    \left.\begin{matrix}\phantom{0}\\ \phantom{0}\end{matrix}\right\}
      {\color{accentwarm}K_1 = \tfrac1{\sqrt2} I}
  \end{array}
$}\par\smallskip
One visible action, scaled the same on both branches: the switch is
fake and the drop is safe, though the temporary ends in $|{+}\rangle$,
not where it started.
\end{minipage}\hfill
\begin{minipage}[t]{0.47\linewidth}
\raggedright
Program (b) sends $v$ to $\tfrac1{\sqrt2}(|0\rangle \otimes v + |1\rangle \otimes Xv)$:
\par\smallskip\centerline{$
  K_{\mathrm{(b)}} \;=\; \tfrac1{\sqrt2}
  \begin{pmatrix}
    {\color{accent}1} & {\color{accent}0} \\
    {\color{accent}0} & {\color{accent}1} \\
    {\color{accentwarm}0} & {\color{accentwarm}1} \\
    {\color{accentwarm}1} & {\color{accentwarm}0}
  \end{pmatrix}
  \begin{array}{@{}l@{}}
    \left.\begin{matrix}\phantom{0}\\ \phantom{0}\end{matrix}\right\}
      {\color{accent}K_0 = \tfrac1{\sqrt2} I} \\
    \left.\begin{matrix}\phantom{0}\\ \phantom{0}\end{matrix}\right\}
      {\color{accentwarm}K_1 = \tfrac1{\sqrt2} X}
  \end{array}
$}\par\smallskip
Two different visible actions. The switch is real, and the drop is
wrong.
\end{minipage}
\end{center}
The checker of \cref{sec:eval} also forms the Gram matrix
$\Gamma_{ij} = \operatorname{tr}(K_i^{\dagger} K_j)$ (\cref{prop:gram}):
$\bigl(\begin{smallmatrix} {\color{accent}1} & {\color{accent}1} \\
{\color{accentwarm}1} & {\color{accentwarm}1} \end{smallmatrix}\bigr)$ of rank
one for (a), the two rows equal, and
$\bigl(\begin{smallmatrix} {\color{accent}1} & {\color{accent}0} \\
{\color{accentwarm}0} & {\color{accentwarm}1} \end{smallmatrix}\bigr)$ of rank
two for (b). The channel is what the
rest of the program sees: after the drop the visible register carries
$\sum_i K_i \rho K_i^{\dagger}$ and nothing else. For program (a) that
is $\rho$, the input untouched. For program (b) it is
$\tfrac12\rho + \tfrac12 X\rho X$, a bit flip with probability one half,
which a continuation that reads $v$ can tell.

When we say \emph{the temporary is left constant} we do not mean
untouched: program (a) moves the temporary to $|{+}\rangle$, for every
input. Nor do we mean \emph{a product at each input}: the copy
$t := t \oplus v$ leaves a product on $|0\rangle$ and on $|1\rangle$
and a Bell state on $|{+}\rangle$, and even a product at \emph{every}
input admits programs the definition excludes, the swap of $t$ and $v$
among them (\cref{prop:productwise}).

\Cref{sec:criterion} says the same thing in the words a language can
use. Program (a) is \emph{weakenable}: weakening is the structural rule that
lets a hypothesis go unused, and program (a) has the denotation of a program
whose temporary is unused, the right-hand picture of \cref{fig:criterion}
with $R = H$ and $W = I$. The equivalence is contextual in the sense
familiar from the $\lambda$-calculus: the contexts are the continuations of
the calculus, a gate on the visible register after the scope followed by a
measurement, and a program is weakenable exactly when no continuation
distinguishes it from some program of that shape (\cref{thm:ctxequiv}). So
\emph{this temporary may be dropped} and \emph{no continuation can tell} are
one statement, and the proportionality test decides it.

\subsection{What the existing checks see}

We encountered three kinds of check in the literature, and each
asks a question about the syntax. A \emph{producer-based}
check---Silq's \code{forget} and Unqomp's synthesised
uncomputation~\cite{Bichsel2020Silq,Paradis2021Unqomp}---asks how the
value was produced: only values produced classically may be forgotten.
A \emph{restore-based} check---Q\# at release and Quipper's
\code{with\_ancilla}
promise~\cite{Svore2018Qsharp,MicrosoftQsharpMemory,Green2013Quipper}---asks
whether the register is back in the state it was allocated in. A
\emph{lifetime} check---Cavy, Qurts and
Qutes~\cite{McNally2021Cavy,Hirata2025Qurts,Faro2026Qutes}---asks who
owns the register and how long it lives, and enforces or synthesises
restoration at the boundary; it answers a different question, whether
cleanup can be inserted, not whether the temporary may be dropped in
the state the program leaves it in. The first two miss different
programs. A producer check refuses a value produced non-classically
that is uncorrelated by the time it is dropped, which is program (a);
a restore check refuses every program leaving the temporary constant
anywhere but $|0\rangle$, however produced. Safety is about
\emph{correlation}, not \emph{location}. Deutsch--Jozsa is refused by
both: its kickback ancilla is prepared in $|{-}\rangle$ by a
non-classical gate, used, and left there (\cref{sec:dj}).

\subsection{Program (c): one text, two verdicts}

\begin{figure}
\centering
\includegraphics{figures/outcome-split}
\caption{Program (c). The allocated state is tilted so that it has a foot in
the measured block and a foot outside it. Reading the program determines the
verdict on each foot: refuse if the block survives, admit if its complement
does. What reading the program cannot determine is which foot survives: the
measurement decides that when the program runs, dynamic lifting hands that one
bit to the compiled code, and the check of \cref{sec:residual} reads it there.
$X$ is the shift $|j\rangle \mapsto |j+1\rangle$ the read applies to the
visible register.}
\Description{The tilted allocated state drawn against the measured block,
with the two outcome branches leading to opposite verdicts on the drop.}
\label{fig:split}
\end{figure}

Program (c) measures the temporary and then reads it, and no later
gate branches on the outcome---yet the verdict does. \Cref{fig:split} is the
whole difficulty in one picture: one program, two reachable outcomes, two
verdicts. Now ask what a
record-oblivious verdict---a function from commands to booleans---can do with
its one boolean. If it says \emph{yes}, it has accepted a drop that is wrong on
the branch of weight $8/9$, so it is unsound. If it says \emph{no}, it has
refused a drop that is safe on the branch of weight $1/9$, so it is incomplete
there. Nothing
about the syntax distinguishes the two branches, and \cref{sec:boundary}
makes that a theorem with a quantifier on it (\cref{sec:boundary}).

Nothing here runs a program. On a measurement-free scope the verdict is a
function of the program text alone; after a measurement the missing argument is
which outcomes survived, not an input. Reading the program fixes the verdict on
every outcome, but not which outcome the run is on, and statistics are the
wrong instrument for the same reason (\cref{sec:channel}).

\subsection{The check at the lift}

The shape of the failure dictates the shape of the fix. A static
analysis must emit its boolean before the run, but the deciding fact
does not exist until the run, after the measurement. So we emit
instead a predicate on the classical outcome, whose table the compiler
computes before the run and consults at the \emph{lift}, where a
measurement outcome becomes available to compiled code. In plain words: the
compiler decides the drop for every outcome the measurement could have, and
the run looks up the one that happened. That outcome is all the run
contributes. For program (c) it is one
classical bit, whether the block $S$ survived, and it selects the row.
\Cref{sec:residual} will prove it sound and complete where every static
analysis is not, at one proportionality test per record.

\begin{keyidea}
A drop is safe when the rest of the program cannot tell: before measurement
the program decides that exactly, by a proportionality test on its Kraus
operators. After measurement
the program is exactly one argument short---the measurement record---and a check at
the lift supplies it.
\end{keyidea}





% Section 3 -- The calculus
\section{The calculus}\label{sec:calculus}

\Cref{sec:overview} was informal. This section fixes the language the theorems are about. Registers are real here and in \cref{sec:criterion}, so $\top$ is the
$\dagger$ of the statements: drop safety is a proportionality condition and
does not see the field. \Cref{sec:gateset}, where phase is the content, is
complex throughout.

\begin{definition}[the affine calculus]\label{def:calculus}
Register and value types, and terms, are
\[
\begin{aligned}
 Q \;&::=\; \tmp[d]\;\mid\;\vis[m],
   \qquad\qquad
   A,B \;::=\; Q\;\mid\;A\otimes B,\\[2pt]
 e \;&::=\; v
   \;\mid\; \langle v_1,v_2\rangle
   \;\mid\; \texttt{let }t=\texttt{new }i_0\texttt{ in }e
   \;\mid\; \texttt{let }(t',v')=U(t,v)\texttt{ in }e
   \;\mid\; \drop\ t\,;e .
\end{aligned}
\]
Contexts are finite partial maps from register names to register types, and
the typing rules are
\begin{ruledisplay}
  \jaxiom{Var}{\jty{v:A}{v}{A}}
\and
  \jrule{Pair}{v_1\neq v_2}
              {\jty{v_1:A,\,v_2:B}{\langle v_1,v_2\rangle}{A\otimes B}}
\and
  \jrule{New}{t\notin\operatorname{dom}(\Delta) \\ i_0\in[d] \\
              \allocFree(e) \\ \jty{\Delta,\,t:\tmp[d]}{e}{C}}
             {\jty{\Delta}{\texttt{let }t=\texttt{new }i_0\texttt{ in }e}{C}}
\and
  \jrule{Gate}{t\neq v \\ t'\neq v' \\
               U:\R^{d}\otimes\R^{m}\to\R^{d}\otimes\R^{m} \\
               U^{\top}U=UU^{\top}=I \\
               \jty{\Delta,\,t':\tmp[d],\,v':\vis[m]}{e}{C}}
              {\jty{\Delta,\,t:\tmp[d],\,v:\vis[m]}
                   {\texttt{let }(t',v')=U(t,v)\texttt{ in }e}{C}}
\and
  \jrule{Drop}{\jty{\Delta}{e}{C}}
              {\jty{\Delta,\,t:\tmp[d]}{\drop\ t\,;e}{C}}
\end{ruledisplay}
\end{definition}

The binders $t'$ and $v'$ bind in $e$, and terms are identified up to
$\alpha$-renaming. A \emph{program} is typed when
$v : \vis[m] \vdash \texttt{let } t = \texttt{new } i_0 \texttt{ in } e : \vis[m]$.
The result type is $\vis[m]$ rather than $\tmp[d]\otimes\vis[m]$,
because we want the type to record that the temporary was weakened
away. A \emph{gate} in \textsc{Gate} always acts on both registers at
once. We call such a gate \emph{joint}, and a gate of the form
$R \otimes I$ or $I \otimes W$ \emph{local}. Contexts are
\emph{affine} in both halves. There is no contraction because copying
a register is not a linear map (no-cloning). There is no
\emph{implicit} weakening because implicit weakening is the entire
subject of the paper: \textsc{Drop} is weakening made explicit, and
under implicit weakening program (b) of \cref{ex:three} would type
with its temporary silently forgotten. The premise $\allocFree(e)$ of
\textsc{New}---the body of an allocation contains no second
allocation---gives a program exactly one temporary scope, and
\cref{sec:many} lifts this.

\textsc{Drop} has no premise about safety, and that is not an
oversight: if $\drop$ demanded safety as a premise, unsafe programs such as
program (b) would not be \emph{expressible}, and the predicates that reject
them could not be compared. \textbf{The result type $\vis[m]$ is a promise, not a
guarantee.} The elaboration of $\drop$ is the identity, the temporary survives as a
latent tensor factor, and the type asserts that no continuation will look at
it. When the drop is unsafe the type is wrong about the denotation, and
detecting exactly that is the rest of the paper. Typing is algorithmic and
types are unique, the outermost constructor of $e$ selecting one rule whose
premises are decided from $\Delta$ and $e$. The calculus admits an arbitrary orthogonal matrix as a gate, and
\cref{sec:gateset} says what a \emph{finite} gate set costs on top of that.

\subsection{The denotation, and drop safety}

The rules say which programs exist. This subsection says what a
program computes, and when its drop is safe. The temporary has state
space $\R^{d}$ and the visible register $\R^{m}$, with $m \ge 1$
throughout. A typed program elaborates to a straight-line list of
gates by recursion on the term: allocation contributes the isometry
$v \mapsto e_{i_0} \otimes v$, a gate contributes its matrix, and
$\drop$ contributes nothing:
\[
\begin{gathered}
  \den{\texttt{let } t = \texttt{new } i_0 \texttt{ in } e}
    = \den{e} \circ (v \mapsto e_{i_0} \otimes v), \qquad
  \den{\texttt{let } (t',v') = U(t,v) \texttt{ in } e}
    = \den{e} \circ U, \\[2pt]
  \den{\drop\ t\,; e} = \den{e}, \qquad \den{v} = I, \qquad
  \den{\langle v_1,v_2\rangle} = \Pi_{v_1 v_2} ,
\end{gathered}
\]
where $\Pi_{v_1v_2}$ permutes the context's tensor factors to put $v_1$ first
and $v_2$ second, and is the identity when they already stand that way.
The \emph{denotation} $K$ of a typed program is this composite, a linear map
$\R^{m} \to \R^{d} \otimes \R^{m}$ computed from the program text alone by
multiplying the matrices of the gates it names. It is an isometry, being an
isometric allocation followed by unitaries. Norm preservation is the
realizability fact typing supplies; absence of leakage is the rank-one
condition below, and the two are different. (\Cref{sec:lifted} adds
measurement, and \cref{thm:meta} there says that this same matrix is what its
machine computes.)

\begin{definition}[Kraus operators, drop safety]\label{def:core}
For a linear map $K : \R^{m} \to \R^{d} \otimes \R^{m}$, its \emph{Kraus
operators} are the $m \times m$ blocks $K_i = (\langle i| \otimes I)\,K$, so
that $Kv = \sum_i |i\rangle \otimes K_i v$. The map $K$ is \emph{drop-safe}
when $K = g \otimes W$ for a vector $g \in \R^{d}$ and a matrix $W$, that is,
$Kv = g \otimes Wv$ for every $v$. A program is drop-safe when its denotation
is.
\end{definition}

The name \emph{Kraus operators} is standard, and the channel is why. Discarding the temporary
leaves the visible register in
\[
  \Phi(\rho) \;=\; \operatorname{Tr}_{\tmp}\bigl(K \rho K^{\top}\bigr)
             \;=\; \sum_i K_i \rho K_i^{\top},
\]
the second equality by expanding the partial trace in the basis $|i\rangle$.
That $K$ is an isometry is exactly $\sum_i
K_i^{\top}K_i = I$, which makes $\Phi$ trace-preserving; over $\C$ every
completely positive map has this form, and every such sum is completely
positive~\cite{Choi1975CP,Kraus1983States}. Read backwards, $K$ is a dilation
of $\Phi$ whose environment is the temporary, and the drop is the partial
trace. The family is not unique---mixing by an isometry, $K'_j = \sum_i
u_{ji}K_i$, gives the same channel---so what belongs to the channel is how
many independent operators it has, not which they are, and that is the number
the criterion reads.

Drop safety says the temporary is \emph{constant} at the drop: it ends in
the one state $g$ whatever the input, unentangled from the visible register.
The force of the definition is the quantifier order, one $g$ for \emph{all}
$v$. Four forms are equivalent, and each is used below: (1) $K = g \otimes
W$; (2) the Kraus operators are pairwise proportional, $K_i = g_i W$; (3)
every $2 \times 2$ minor of the $d \times m^{2}$ matrix $(i, (b,c)) \mapsto
K_{(i,b),c}$ vanishes; (4) the Gram matrix $\Gamma_{ij} = \operatorname{tr}(K_i^{\top}
K_j)$ of the Kraus operators has rank at most one. Form (3) is a finite
computation on exact numbers, valid over any field, and \cref{sec:cost}
prices it. Form (4) is form (3) compressed to $d \times d$: $\Gamma = MM^{\top}$ for
the matrix $M$ of form (3), a Gram matrix has the rank of the family it is
built from, and it is the form the checker cross-checks against. Since $K$ is
an isometry,
$\operatorname{tr}\Gamma = \sum_i \operatorname{tr}(K_i^{\top}K_i) = m$, and
a positive semidefinite matrix of trace $m$ has rank at most one exactly when
its eigenvalues are $m$ once and $0$ otherwise, that is, when
\[
  \Gamma^{2} = m\,\Gamma,
  \qquad\text{equivalently}\qquad
  \Bigl(\frac{\Gamma}{m}\Bigr)^{2} = \frac{\Gamma}{m} :
\]
the drop is safe exactly when $\Gamma/m$ is \emph{idempotent}, a rank-one
projector, and then it is $gg^{\top}$ for the unit final state $g$, since
$\operatorname{tr}(W^{\top}W) = m$ for the isometry $W$. That projector is
the temporary's own state at the maximally mixed visible input
(Appendix~A.5), and a state is idempotent exactly when it is pure. So the
criterion has a one-line reading on the temporary alone: \emph{a drop is safe
exactly when the temporary comes out pure at the maximally mixed visible
input}. That reference input is not a convenience, as the next proposition
shows. The pair $(g,W)$ is determined up to $(\lambda g, \lambda^{-1}W)$, so
``unit $g$'' is a choice of representative; the zero map is drop-safe.

Here is why that single reference input is enough, and why the obvious
alternative is not the same test. A reader meeting \cref{def:core} may well
ask what is wrong with checking the joint state \emph{at each input
separately}: run the program on $v$, see whether what comes out is a product
of a temporary state with a visible state, and repeat for every $v$. That is a
weaker demand---it lets the temporary's state depend on $v$, where drop safety
does not---and it is weak enough to admit programs that leak. The next
proposition says exactly how much weaker: pointwise productness has only two
solutions, and only the first of them is a safe drop.

\begin{revisingblock}
\stmtProductwise
\end{revisingblock}

The proof (Appendix~A.1) is linearity: two inputs whose temporary factors are
not parallel force, through $K(v_1 + v_2)$, their visible factors to be
parallel---so one of the two factors has to stand still, and which one it is
gives the two cases. The second case is the leak. It puts the input into the
temporary and leaves the visible register alone, and the plainest instance is
the \textbf{swap} of $t$ and $v$ at $d = m$: it is a product at every input,
its Kraus operators are $e_{i_0}e_i^{\top}$, and unsafe for every $m \ge
2$, having moved the whole input into the register about to be discarded. So
\emph{a product at every input} admits programs \cref{def:core} excludes, and
the quantifier order---one $g$ for \emph{all} $v$---rules them out.

\stmtGram

The proof (Appendix~A.1) is short. Part (1) is the Knill--Laflamme
condition~\cite{Knill1997Theory} with the whole visible register as the
code, the $K_i$ as the errors and the temporary as the environment, which
therefore learns nothing~\cite{Kretschmann2008Tradeoff}. Its second half is
noninterference~\cite{Goguen1982Security,Ying2013Noninterference}, over $\C$
only---a real rotation refutes it (Appendix~A.1)---and the impurity $\iota$ of
\cref{sec:cost} is its quantitative degree. Part
(2) reads Choi's rank~\cite{Choi1975CP} as a resource, the least width a
temporary needs for the observed visible behaviour. Part (3) makes $\Gamma$
a certificate: nothing the continuation does changes it.


One lemma about that family is used twice below, once for the criterion of
\cref{sec:criterion} and once for its reading as a channel, so it is stated
here where the Kraus operators are introduced.

\stmtProp

The proof (Appendix~A.1) reads the hypothesis one matrix entry at a time, so
it is worth writing the two sides out. Both are $m \times m$ and symmetric;
at $m = 2$, abbreviating $u_i = K_i x$ and $z = Mx$, the hypothesis says
\[
  \sum_i
  \begin{pmatrix}
    {\color{accent}(u_i)_0^{2}} & (u_i)_0 (u_i)_1 \\
    (u_i)_1 (u_i)_0 & {\color{accentwarm}(u_i)_1^{2}}
  \end{pmatrix}
  \;=\;
  \begin{pmatrix}
    {\color{accent}z_0^{2}} & z_0 z_1 \\
    z_1 z_0 & {\color{accentwarm}z_1^{2}}
  \end{pmatrix} .
\]
The coloured diagonal settles each row on its own. Where $z_b = 0$ a sum of
squares is zero, so every $(u_i)_b$ vanishes with it; where $z_b \neq 0$ the
ratios $(u_i)_b / z_b$ form a unit vector $c^{(b)}$ indexed by $i$. That is
the sense in which each row of each $K_i$ is a multiple of the corresponding
row of $M$. The uncoloured off-diagonal stops the multiples drifting from row
to row: it reads $\langle c^{(0)}, c^{(1)}\rangle = 1$, and two unit vectors
with inner product one are equal, so the rows share one $c$: the statement's
$c_i$.

\subsection{The constructs in existing systems}

Every construct in the calculus is one that existing systems already have.
$\texttt{new}$ is Silq's local quantum variable, Q\#'s \code{use}, and
Quipper's \code{with\_ancilla}; Q\#'s \code{borrow} is the dirty variant,
which carries no initialisation promise. $\drop$ is Silq's
\code{forget}, the end of a Q\# \code{use} block, and the close of
\code{with\_ancilla}. The side condition is Silq's \code{qfree} producer,
checked by its type checker, Q\#'s requirement that a qubit be released in
the state it was allocated in, checked at run time on simulators, and in Quipper the programmer's promise,
checked by nothing. The $\lift(\meas)$ of \cref{sec:lifted} is Q\#'s \code{M}
followed by classical control, and Quipper's dynamic lifting. The one mismatch
runs the other way:
Silq's \code{measure} consumes the register, where a lift leaves the temporary
live and restricted to the surviving block.
\Cref{sec:eval} says which of these are transcripts and which documentation.

% Section 4 -- The criterion
\section{The criterion}\label{sec:criterion}

Both of the paper's negatives are stated against a positive, and this is
it. For the measurement-free calculus we can say exactly when $\drop$ is
safe, and say it not as ``the denotation factors''---\cref{def:core},
quantified over matrices---but as a statement quantified over
\emph{programs}.

\begin{definition}[weakening witness]\label{def:sw}
Write $\sep(R,W)$ for the term
\[
  \texttt{let } (t_1,v_1) = (R\otimes 1)(t,v) \texttt{ in }
  \texttt{let } (t_2,v_2) = (1\otimes W)(t_1,v_1) \texttt{ in }
  \drop\ t_2\,;\ v_2
\]
--- \emph{prepare the temporary once, compute on the visible register, drop}.
A program $P$ is \emph{weakenable} when some \emph{typed} program of that
shape, with the same allocation, has the same denotation.
\end{definition}

Two readings of that last sentence have to be kept apart, and the section's
whole content sits in the gap between them.

\begin{description}
\item[\emph{being} $\sep(R,W)$] is a fact about the program \emph{text}: this
  text has one gate on the temporary, then one on the visible register, and
  nothing joint anywhere.
\item[\emph{being weakenable}] is a fact about the \emph{denotation}:
  somewhere in the calculus there is a text of that shape computing the same
  function---whether or not our program looks remotely like it.
\end{description}

The second is by far the weaker demand, and Deutsch--Jozsa is the case to hold
on to. Its text has the joint gate $U_f$ sitting in the middle of it and
resembles $\sep(R,W)$ not at all; yet its denotation is $|{-}\rangle \otimes
HD_fHv$ (\cref{sec:dj} works this out), and the separable program that
prepares $|{-}\rangle$ and then applies $HD_fH$ has exactly that denotation.
The algorithm is weakenable without being separable, and by
\cref{thm:ctxequiv} no continuation will tell the two apart. Of the
artifact's seventeen catalogue programs \textbf{two} are literally of the
shape; \textbf{twelve} are safe and so, by the next theorem, have a witness
they need not resemble. Appendix~B.1 of the supplement prints three of them
beside their programs.

\begin{theorem}[the criterion]\label{thm:criterion}
A typed program is weakenable if and only if it is drop-safe.
\end{theorem}
\begin{proof}\leavevmode
\begin{steps}
\item ($\Rightarrow$) A witness $\sep(R,W)$ sends $e_{i_0} \otimes v \mapsto
  Re_{i_0} \otimes v \mapsto Re_{i_0} \otimes Wv$, which is $g \otimes W$
  with $g = Re_{i_0}$, and $\drop$ changes nothing.
\item ($\Leftarrow$) Let $P$ be typed with denotation $Kv = g \otimes Wv$.
  Since $K$ is an isometry, $\|g\|\,\|Wv\| = \|v\|$ for every $v$. Rescale
  to $\|g\| = 1$, and $W$ is an isometry of $\R^{m}$, hence orthogonal,
  hence a legal gate.
\item It remains to prepare $g$ from $e_{i_0}$ by one gate. If $g = e_{i_0}$
  take $R = I$. Otherwise put $w = e_{i_0} - g$ and let $R = I -
  2ww^{\top}/\|w\|^{2}$, the reflection through the hyperplane $w^{\perp}$:
  it is orthogonal, and since $\langle w, e_{i_0}\rangle = 1 - \langle g,
  e_{i_0}\rangle = \tfrac12\|w\|^{2}$ it sends $e_{i_0}$ to $e_{i_0} - w =
  g$. Then $\sep(R,W)$ is typed, has the same allocation, and denotes
  $K$.\qedhere
\end{steps}
\end{proof}

Step $\langle 3\rangle$ is the substance: every unit final state is
preparable by one reflection or the identity, hence one \textsc{Gate}, so the certificate
$(g,W)$ is realised by a program, the step a finite gate set threatens
(\cref{sec:gateset}). Step $\langle 2\rangle$ is the whole contribution of
the type system: its gates are unitary, so $K$ is an isometry. Without
typing that must be assumed: for arbitrary linear gates, weakenability is
drop safety together with $\|Kv\| = \|v\|$ for every $v$, and the clause
cannot be dropped, since at $d = m = 1$ the map that doubles the visible
register is drop-safe and not weakenable (Appendix~A.1). What is real-only
here is the reflection witness; over $\C$ the factors are determined up to a
compensating phase.

\subsection{The criterion is contextual indistinguishability}

\Cref{def:sw} compares a program with a witness at every input. That is
not the same as comparing them in every \emph{context}. A
\emph{continuation} appends one visible gate $C$ after the scope's last gate,
drops the temporary, and reads the visible register in the standard basis.
What it observes at input $x$ is the distribution $b \mapsto \sum_i \langle
e_b, C K_i x\rangle^{2}$: the Born statistics of the visible register,
temporary traced out. A class $F$ of visible gates \emph{distinguishes} two
programs when some $C \in F$ and some input give different distributions, and
a \emph{separable program} is a typed program of the shape $\sep(R,W)$.

\begin{revisingblock}
\stmtCtxEquiv
\end{revisingblock}

The proof (Appendix~A.1) reads the Born statistics after each reflection:
agreement at every reflection says the symmetric matrices $\sum_i (K_i
x)(K_i x)^{\top}$ and $(Wx)(Wx)^{\top}$ have the same quadratic form, so they
coincide, and \cref{lem:prop} gives $K_i = c_i W$.

Three remarks. A larger class distinguishes more, so the theorem holds
for all visible gates and for continuations that also touch the temporary.
The class cannot be $\{I\}$: the unsafe copy $t := t \oplus v$ has the same standard-basis statistics as doing nothing at every input, since $\sum_i \langle e_b, P_i x\rangle^{2} =
x_b^{2}$ for its Kraus projectors $P_i$. The exact minimal class is open. And
the theorem is about programs whose gates keep their register names, since the
gate rule may rename its output binders.

\subsection{Worked: Deutsch--Jozsa}\label{sec:dj}

The criterion on the textbook's first algorithm, where it says what a producer-based check cannot:
\begin{prog}
let t = new |0> in t := X t; t := H t; & -- the ancilla is prepared in |->\\
v := H v; v,t := Uf(v,t); v := H v; & -- phase kickback\\
drop t; v & -- the drop under discussion
\end{prog}
Take $f : \{0,1\} \to \{0,1\}$, so that $U_f|j,t\rangle = |j, t \oplus
f(j)\rangle$ and $D_f = \operatorname{diag}((-1)^{f(0)}, (-1)^{f(1)})$. At an \emph{arbitrary} input
$\sum_j c_j|j\rangle$ to $U_f$---not only the intended one---the state after
it is $\bigl(\sum_j (-1)^{f(j)} c_j |j\rangle\bigr) \otimes |{-}\rangle$.
The temporary factor is $|{-}\rangle$ whatever the input, so the scope's
denotation is $g \otimes Wv$ with $g = |{-}\rangle$ and $W = HD_fH$, which is
$\pm I$ for constant $f$ and $\pm X$ for balanced $f$. By
\cref{thm:criterion} the drop is safe, and by \cref{thm:ctxequiv} no
continuation can distinguish this program from a separable one. It is
program (a) with an algorithm attached, and every rule that reads the program
text refuses it, for one of the two reasons of \cref{sec:overview}. The
producer-based rules refuse it because $H$ is not \code{qfree}: Silq's
automatic \code{forget} fails in our transcripts with \emph{cannot synthesize
forget expression}, and Unqomp's synthesis is defeated by the same
gate~\cite{Bichsel2020Silq,Paradis2021Unqomp}. The restore-based rules refuse
it because $|{-}\rangle$ is not $|0\rangle$: Q\#'s release
requirement~\cite{Svore2018Qsharp,MicrosoftQsharpMemory}, the promise
Quipper's \code{with\_ancilla} closes on~\cite{Green2013Quipper}, Qutes'
Reclamation~\cite[Def.~6]{Faro2026Qutes}, and the $\Pi$ family's fractional
reuse rule~\cite{Chen2020Fractional}. None of these four turns on $f$: none of
them looks at the oracle, and each refuses on the state the temporary is left
on. \code{dqt pool} runs the last of them. The verification logics of
\cref{sec:related} can admit the drop, but from a separability assertion the
programmer supplies and proves, not from a check on the text.

\subsection{What the pool receives, and the rule it already had}

\begin{planned}{the certificate, and the rule the classical lineage already had}
What the free pool receives when a drop is admitted, and what a record-indexed
pool would have to receive once allocation may depend on an outcome; and how
the $\Pi$ family's fractional reuse rule~\cite{Chen2020Fractional} reads
against the criterion. Being written; nothing here is claimed.
\end{planned}


% Section 5 -- What restricting the gates costs
\section{What restricting the gates costs}\label{sec:gateset}

\Cref{thm:criterion}'s right-to-left direction used that \emph{every}
unit state is preparable in one gate, because \cref{def:calculus} admits an
arbitrary orthogonal $U$. A compiler emits words over a fixed finite
repertoire, and there preparability is a real constraint: the biconditional
can acquire a side condition. We state it at the \emph{ring} the
repertoire's entries generate: every word over the repertoire is a gate over
that ring, so an obstruction over the ring is one for the repertoire. The condition bites over some rings, and
not over Clifford$+T$.

Fix a \emph{gate ring} $\mathcal A$: a subring of $\C$ closed under
conjugation, the ring the gate entries generate. A gate over
$\mathcal A$ is an $\mathcal A$-valued unitary. Transport back to a
concrete repertoire is not automatic: it needs exact synthesis
\emph{within the witness's allocation budget}, and the multiqubit
results permit a clean
ancilla~\cite{Kliuchnikov2013ExactSynthesis,Giles2013ExactSynthesis,Amy2020Restricted},
so we infer no finite-repertoire corollary. Say $\mathcal A$ is
\emph{completable} in dimension $d$ when every unit vector in
$\mathcal A^{d}$ is a column of a unitary over $\mathcal A$.
Throughout this section the temporary has dimension $d \ge 2$ and is
allocated at $|0\rangle$, the visible register is arbitrary, and
scalars are complex. When a program $U$ is drop-safe its final state
$g$ is determined up to phase, and the object of interest is the
rank-one projector $gg^{\dagger}$, computed directly from the entries
of $U$ without first finding $g$:
\[
  \mathrm{proj}(U)_{i\ell} \;=\; \sum_{{\color{accentplum}c}}
     U_{({\color{accent}i},\,0),\,({\color{accentwarm}0},\,{\color{accentplum}c})}\;
     \overline{U_{({\color{accent}\ell},\,0),\,({\color{accentwarm}0},\,{\color{accentplum}c})}} .
\]
Read the four slots of $U_{(i,b),(j,c)}$ as \emph{amplitude from temporary
$j$, visible $c$, to temporary $i$, visible $b$}, and the colours say which is
which. {\color{accentwarm}Warm} is the allocation: the input temporary is
pinned to $0$, because that is where a program starts it.
{\color{accentplum}Plum} is the visible index, and it is
\emph{summed}---nothing here depends on which visible basis vector went in,
which is why no clause about the visible register appears in the theorem
below. {\color{accent}Blue} is what the pinnings keep: the picked cells, and the two
output temporary levels $i$ and $\ell$, which are the free indices---the row
and column of the $d \times d$ matrix being built. The output visible index is pinned to $0$ in both factors, and for a
drop-safe $U$ any other choice gives the same matrix: the entries are then
$U_{(i,b),(0,c)} = g_i W_{bc}$, so the sum is $g_i \overline{g_\ell}$ times the
squared norm of one row of the isometry $W$, which is $1$.
\Cref{fig:proj-index} is that formula drawn: the two pinnings are a block
column and a row inside it, and what survives is a slab against its own
conjugate.

\begin{figure}[!ht]
\centering
\includegraphics{figures/proj-index}
\Description{A six by six matrix drawn as a three by three array of two by
two blocks. The first block column is shaded, and the first row of each block
in it is picked out. Those picked cells form a three by two slab, and the
slab times its own conjugate transpose is a three by three matrix.}
\caption{\textbf{Look at the block column}: the allocation pins the input
temporary to $0$, and inside that column the output visible index is pinned to
$0$ as well. What is left is a $d \times m$ slab, and $\mathrm{proj}(U)$ is
that slab against its own conjugate with the visible index summed away. Drawn
at $d = 3$, $m = 2$.}
\label{fig:proj-index}
\end{figure}

Say $\mathcal A$ \emph{splits} a rank-one projection $P$ when $P =
zz^{\dagger}$ for a vector $z$ over $\mathcal A$. The two halves of that
sentence are different demands, and \cref{fig:splitting} keeps them apart:
$P$ may have every entry in $\mathcal A$ while no vector over $\mathcal A$
splits it, because the splitting vector's entries are square roots of $P$'s.

\begin{figure}[!ht]
\centering
\includegraphics{figures/splitting}
\Description{On the left, a d by d grid of cells with a column vector to its
left and a conjugated row vector above it, showing that entry i, l of the
matrix is the i-th entry of the vector times the conjugate of its l-th. On the
right, the projection onto plus, over two rings, with the splitting vector
available in one and not the other.}
\caption{\textbf{Look at the two rows on the right}: the projector is the
same matrix in both, and only the ring changes. Splitting asks for a vector
over the ring, which is strictly more than asking that the matrix's entries
lie in it.}
\label{fig:splitting}
\end{figure}

\stmtRelativised

\Cref{fig:ring-repertoire} is this section in one path: repertoire to ring,
ring to the splitting test, and the crossed way back.

\begin{figure}[!ht]
\centering
\includegraphics{figures/ring-repertoire}
\Description{A flowchart from a finite repertoire to its coefficient ring and
then to a splitting test. Success yields a separable program; failure yields a
program that is not weakenable. A crossed dashed arrow shows that
ring-valued separation does not by itself synthesise a word over the finite
repertoire.}
\caption{\revisinginline{Which way section~\ref{sec:gateset} runs.
\textbf{Look at the branch}: splitting is a side condition the ring does not
supply, and \cref{thm:relativised} answers it both ways. \textbf{Then the
dashed arrow}: the way back to words is exact synthesis, and not
free---the case that fails is below the theorem.}}
\label{fig:ring-repertoire}
\end{figure}


The proof (Appendix~A.2) reads $\mathrm{proj}(U) = gg^{\dagger}$ off a
drop-safe $U$. A witness over $\mathcal A$ prepares $g$ by a gate over
$\mathcal A$, so the projector splits, and conversely a splitting vector $z
= \lambda g$ is completed to a gate by completability, while the visible
factor is recovered over $\mathcal A$ by contracting $U$ against $\bar z$.

Three things in the statement are not the expected shape. \emph{The
projector is free}: $\mathrm{proj}(U)$ is a polynomial in the entries of $U$
and their conjugates, so it lies over $\mathcal A$, and the side condition
is a rank-one factorisation, not a matrix square root. \emph{No clause for
the visible factor is needed}, and \emph{the visible register plays no
role}, whatever its width. Every gate ring is completable at $d = 2$, and five rings are
completable in every dimension, by the column lemmas of the exact-synthesis
literature: $\mathcal R = \Z[1/\sqrt2, i]$, the ring of
Clifford$+T$~\cite[Lemma~5]{Giles2013ExactSynthesis}, and $\Z[1/2]$,
$\Z[1/\sqrt2]$, $\Z[i/\sqrt2]$ and $\Z[1/2, i]$, its integral, real,
imaginary and Gaussian restrictions~\cite[Lemmas~5.4, 5.14, 5.19,
5.25]{Amy2020Restricted}.

\stmtCliffordT

The ring $\mathcal R$ is exactly the one in which single-qubit
Clifford$+T$ operators live~\cite{Kliuchnikov2013ExactSynthesis}: a $2 \times
2$ unitary over $\mathcal R$ is an ancilla-free Clifford$+T$ word up to a
global phase, and conversely. That reading needs the \emph{whole} unitary to be one qubit,
temporary and visible together---not merely the visible factor---and a
program with a temporary is already past it. Beyond it a gate over $\mathcal
R$ is a macro-gate, and synthesising one in Clifford$+T$ can need a clean
ancilla~\cite{Giles2013ExactSynthesis}, which a witness reusing the same
allocation does not have. Controlled-$T$ is such a case: every ancilla-free
two-qubit word in $H$, $T$ and CNOT has determinant a fourth root of unity,
$\det$ controlled-$T = e^{i\pi/4}$, and the ring's scalars move a two-qubit
determinant only by $\pm1$. The theorem is about the ring, and we do not
claim it for a finite repertoire. Part (1) is proved in Appendix~A.2 from one
classical fact, that $\Z[e^{i\pi/4}]$ has class
number one~\cite[Ch.~11]{Washington1997Cyclotomic}: a nonzero column of the
projection, divided by the greatest common divisor of its entries, has for
its squared norm a unit of $\Z[1/\sqrt2]$ that is a sum of norms, hence
totally positive, hence $|\varepsilon|^{2}$ for a unit $\varepsilon$ of
$\mathcal R$, and that column divided by $\varepsilon$ is $z$. Part (2) is
part (1) inside \cref{thm:relativised}, the column lemma supplying
completability.

One ring below, the side condition is real. Over $\Z[1/\sqrt2]$---the ring of
the real Clifford$+T$ gates---the projector of the $H$-type magic state
$\cos\tfrac\pi8|0\rangle + \sin\tfrac\pi8|1\rangle$ has entries in the ring
and does not split, while over $\mathcal R$ it does; and a single gate
$\mathit{rot} \otimes \mathit{rot}$, for $\mathit{rot}$ the rotation of a
qubit by $\pi/8$, realises the separation as a program, drop-safe with that
magic state as its final state and weakenable over $\mathcal R$ but not over
$\Z[1/\sqrt2]$. Appendix~A.2 computes both. What changed is not the final
state but the ring available to name it.

\stmtRings

The proof (Appendix~A.2) is \cref{thm:cliffordT}'s argument again: the two complex rings are principal ideal domains
whose totally positive real units are norms of units, while over a real
ring a splitting vector can only be $\pm g$. Of the five rings named here,
the three complex ones split every projection and the two real ones each
refuse a named projection.

\subsection{What deciding costs}\label{sec:cost}
The criterion of \cref{def:core} is a finite computation on the denotation,
and the denotation is exponentially large in the width. This subsection prices
it. Scalars are complex here, so $\Gamma_{ij} = \operatorname{tr}(K_i^{\dagger}
K_j)$. Three definitions are needed first, and they are the only ones.

\begin{description}
\item[the input, a \emph{circuit program}] has qubit registers, so $d = 2^{k}$
  and $m = 2^{n}$, with the temporary allocated at $|0^{k}\rangle$ and gates on
  at most three qubits each. A gate is given by its matrix, each entry an
  algebraic number presented exactly---its minimal polynomial, and an isolating
  rectangle in binary. The program's \emph{size} counts the gates, the qubits
  and that presentation.
\item[the leak, $\iota(K)$] is measured on the temporary. Feed the visible
  register an input $\psi$, possibly entangled with a reference register, and
  let $\rho_t(\psi)$ be the temporary's state after the run. Its
  \emph{impurity} $1 - \operatorname{tr}\rho_t(\psi)^{2}$ vanishes exactly
  when that state is pure, and $\iota(K) = \sup_\psi\bigl(1 -
  \operatorname{tr}\rho_t(\psi)^{2}\bigr)$ is the worst case over inputs.
\item[the leak at one input, $\bar\iota(K)$] is the same number read at a
  single $\psi$, one whose visible marginal is maximally mixed. That is well
  defined because $\rho_t(\psi)$ depends on the visible marginal alone, and it
  is what makes the quantity computable rather than a supremum.
\end{description}

At an input whose visible marginal is maximally mixed---the one $\bar\iota$
is measured at---the temporary is in the state $\Gamma^{\top}/m$, by the
entrywise identity $\operatorname{Tr}_v[KK^{\dagger}] = \Gamma^{\top}$
(Appendix~A.5); so $\bar\iota(K) = 1 - \operatorname{tr}(\Gamma^{2})/m^{2}$,
and the drop is safe exactly when $\bar\iota(K) = 0$, exactly when $\iota(K)
= 0$. That is form (4) read physically, and it is what the promise below is
stated on.

\stmtClassicalCost

Membership is two basis inputs that land the temporary in different
basis states. Hardness is Bennett's compute--copy--uncompute~\cite{Bennett1973Logical}
with the target counted into the temporary: a program that evaluates a
Boolean circuit $B$, copies $B(x)$ into a target and uncomputes its scratch
drops the two together safely exactly when $B$ is constant. Constancy is not
unsatisfiability, so the reduction emits a fixed unsafe program when
$B(0^{n}) = 1$; what it emits is drop-safe exactly when $B$ is unsatisfiable. Liu et al.\ show the same bound for a
different predicate, whether some unitary can restore the
ancilla~\cite{Liu2026Uncomputation}, and that program is restorable for every
$B$ (Appendix~A.5).

\begin{revisingblock}
\stmtAverageCost

The algorithm: run the program twice on independent uniformly random
basis inputs and swap-test the two temporaries, which accepts
with probability $1 - \bar\iota(K)/2$~\cite{Buhrman2001Fingerprinting}. The two readings of the leak differ.
For $m \ge 2$ Grover's reflection~\cite{Grover1996Search} $U = I -
2|w\rangle\langle w|$ has nonzero $+1$ and $-1$ eigenspaces, and the program that puts
the temporary in $|{+}\rangle$, applies $U$ controlled by it and drops it has
$\iota = \tfrac12$, the most a qubit can carry, but
$\operatorname{tr}(U)/m = 1 - 2/m$, so $\bar\iota = 2/m - 2/m^{2}$,
and the swap test needs order $m$ repetitions to see it,
exponentially many in the width. Without a promise gap the verdict is
coNP-hard over any gate set containing $X$, CNOT and Toffoli, and lies in
PSPACE when the gate entries come from a fixed number field
(Appendix~A.5).
\end{revisingblock}

\begin{planned}{the price on a restricted gate set}
On which gate sets is the verdict cheap, and what does a gate one step outside
such a set cost? Being written; nothing here is claimed.
\end{planned}


% Section 6 -- Measurement, and the lifted calculus
\section{Measurement, and the lifted calculus}\label{sec:lifted}

Program (c) needs a language in which a measured outcome becomes a
classical parameter. That is a second language, and from here the two are
kept apart by name: a \emph{program} is a typed term of \cref{def:calculus},
which is what every section up to here was about, and a \emph{command} is a
typed $c$ of the calculus just below---which is what (c) becomes on that
reading. The rest of
the paper is about commands.

The second calculus is two-sorted, in the style of Proto-Quipper's
\emph{dynamic lifting}~\cite{Fu2023DynamicLifting}---the Quipper term for
reading a measurement outcome into classical control:
\[
\begin{aligned}
  c ::={}& \texttt{halt} \mid \texttt{gate}\ U\,;c \mid \texttt{vgate}\ V\,;c
         \mid \drop\ t\,;c \\
      \mid{}& \texttt{let } p = \lift(\meas\ t \in S)\texttt{ in } c
         \mid \texttt{if } e \texttt{ then } c_1 \texttt{ else } c_2
\end{aligned}
\]
with two contexts: parameters are classical, duplicable and weakenable, and the
quantum state is affine, forgettable only by $\drop$, after which nothing
may name it. There is no allocation form, and none is missing: a command is
typed at a flag that starts the temporary live, and the machine starts it at
$e_{i_0}$, so the allocation sits in the judgement rather than in the syntax.

That is what $\allocFree$ does for \cref{def:calculus}, said the other way
round---one temporary, allocated once, before anything runs. One rule
connects the two sorts: a lift on a live temporary
measures whether it lies in the block $S$, keeps the surviving branch without
renormalising, and binds the answer as a parameter the continuation may test.
The \emph{measurement record}, the word of the quantum-trajectory literature~\cite{Wiseman2010Measurement}, is the sequence of which outcome survived at each
$\lift$, supplied from outside the program. Branch
nondeterminism is recovered by quantifying over records.

What a command \emph{does} still has to be said, and the lift is the reason to
say it with a machine rather than a function. A lift has two futures
and the text does not choose between them; only the record does. So the
machine is handed its record from outside, as a tape it consumes one entry per
lift; everything else is then determined. A
\emph{configuration} $\langle c;\psi;\eta;\sigma\rangle$ carries the command
still to run, the joint state $\psi$, the store $\eta$ of parameters bound so
far, and the record's unread suffix $\sigma$.

Two abbreviations make the lift rule readable. A lift on the block $S$ keeps
either the part of the temporary inside $S$ or the part outside it, so write
\[
  P_{\mathsf{in}} = P_S, \qquad P_{\mathsf{out}} = I - P_S,
  \qquad \beta(\mathsf{in}) = 1, \qquad \beta(\mathsf{out}) = 0
\]
for those two projections and for the boolean each of them reports back to the
program. Three reductions then carry the weight; a visible gate and a
conditional do the obvious thing, and Appendix~A.1 has them.
\[
\begin{aligned}
 \langle \texttt{gate}\ U\,;c;\psi;\eta;\sigma\rangle
 &\longrightarrow
 \langle c;U\psi;\eta;\sigma\rangle,\\
 \langle \drop\ t;c;\psi;\eta;\sigma\rangle
 &\longrightarrow
 \langle c;\psi;\eta;\sigma\rangle,\\
 \langle \texttt{let }p=\lift(\meas\ t\in S)\texttt{ in }c;
          \psi;\eta;b\cdot\sigma\rangle
 &\longrightarrow
 \langle c;(P_b\otimes I)\psi;\eta[p\mapsto\beta(b)];\sigma\rangle.
\end{aligned}
\]
Read them in order. A gate applies its matrix and moves nothing else.
A drop applies \emph{nothing}, which is exactly the point: the temporary is
still sitting in $\psi$, and all that has happened is that the typing has
stopped anyone naming it. A lift takes the next entry $b$ off the record,
projects the temporary onto that half of the block, does \emph{not}
renormalise---so the weight the branch keeps is carried rather than thrown
away---and binds $\beta(b)$ into the store for a later conditional to test.

Because the record was fixed before the run, all of this multiplies out into
one matrix. Write $E_o(c)$, the \emph{selected denotation} of $c$ along $o$,
for the composite of the gates and block projections the run performs, read
off $c$ by a recursion that Appendix~A.1 gives together with the typing rules
and the remaining reductions; \cref{thm:meta}, just below, is the statement
that this matrix and the machine agree. A record is \emph{reachable} for $c$
when $E_o(c) \ne 0$. The drop is \emph{safe under $o$} when $E_o(c)$ is drop-safe, so an
unreachable record is safe. The \emph{weight} of the run at a unit input
$v$ is $\|E_o(c)v\|^{2}$.

\stmtMeta

The proof, a pair of inductions over commands, is in Appendix~A.1. What the
theorem buys is that every later statement about a command is a statement
about the single linear map $E_o(c)$---and that map is of the kind
\cref{def:core} is about. The machine starts at $e_{i_0} \otimes v$, so
$E_o(c)(v) = E_{o,\emptyset}(c)(e_{i_0}\otimes v)$: an operator on
$\R^{d}\otimes\R^{m}$ precomposed with the allocation isometry, hence a map
$\R^{m} \to \R^{d}\otimes\R^{m}$, which is the type of $K$. Kraus operators,
drop safety and the Gram matrix are therefore read off $E_o(c)$ exactly as
off a program's $K$, one record at a time. Only the normalisation moves: a
program's $K$ is an isometry, so $\operatorname{tr}\Gamma = m$, while
$E_o(c)$ carries the record's block projections and
$\operatorname{tr}\Gamma$ is instead the weight that record keeps, summed
over the visible basis. \Cref{thm:bridge} is that bookkeeping.

The two calculi meet at their denotations, not at their grammars. A typed
program has no lift, and the gate list its elaboration produces \emph{is} a
lift-free command with the same denotation; the lift-free commands are
slightly the larger class, since one may apply a visible gate after the drop
and no program can write that. So measurement does not break the matrix
picture. It indexes it.

\subsection{The same verdict, as a channel}\label{sec:channel}

\Cref{thm:meta} paves the way for reasoning with channels: along a record a
command is one linear map $E_o(c)$, so the verdict can be stated on the
channel that map induces. Fix a measurement record $o$, write $V_o = E_o(c)$,
and build the discard channel of \cref{sec:calculus} from $V_o$ instead of
from a program's $K$: $\Phi_o(\rho) = \sum_i K_i \rho K_i^{\top}$ with $K_i$
the Kraus operators of $V_o$. Because $V_o$ carries the block projections of the record
it is not an isometry, so $\Phi_o$ is trace-\emph{non}increasing rather than
trace-preserving---a \emph{subchannel}, whose surviving trace is the weight
of the branch. The channel is $\sum_o \Phi_o$; the subscript is load-bearing, because
without it there is no single branch to speak of. The proportionality
lemma of \cref{sec:calculus} now does its second job.


\begin{theorem}[drop safety is Kraus rank one]\label{thm:bridge}
Fix a measurement record $o$. The drop is safe under $o$ if and only if the
subchannel $\Phi_o$ has a representation with a single Kraus operator. At a unit input $v$ the
weight of the run is $\operatorname{tr} \Phi_o(v v^{\top})$.
\end{theorem}
\begin{proof}
$\operatorname{tr} \Phi_o(vv^{\top}) = \sum_i \|K_i v\|^{2} = \|V_o v\|^{2}$
is the weight. If the drop is safe, $K_i = c_i W$ and the sum collapses to
$\Phi_o(\rho) = (cW) \rho (cW)^{\top}$ with $c^{2} = \sum_i c_i^{2}$. Conversely, if
$\Phi_o(\rho) = M \rho M^{\top}$ for every $\rho$, then at $\rho = xx^{\top}$ the
hypothesis of \cref{lem:prop} holds, so $K_i = c_i M$ and $V_o = (c_i)_i
\otimes M$ is drop-safe. Nothing assumed $V_o$ isometric, and the zero map has the
single Kraus operator $0$, so an unreachable record is safe here too.
\end{proof}

One caution, and it is why the paper is indexed by measurement records
rather than by channels. Two well-typed commands on a temporary qubit and a
visible qubit can have the same summed channel $\Phi_0 + \Phi_1$ and the
same weight on every branch at every input, with opposite verdicts on every
branch: $H$ then $\mathrm{CZ}$ then a lift, against a lift then a controlled
write, both summing to the dephasing channel (Appendix~A.1). So an analysis
that reads the channel a program implements, or samples its statistics,
cannot decide the drop. One that reads the program can, and its instance is
the calculus above.


% Section 7 -- The boundary of one verdict per command
\section{The boundary of one verdict per command}\label{sec:boundary}

\begin{planned}{the boundary, and where it starts}
What a verdict that does not read the measurement record must give up, and
where that boundary starts. Not an undecidability result: safety under every
reachable record is decided by running \cref{thm:criterion}'s rank test once
per record. Being written; nothing here is claimed.
\end{planned}


% Section 8 -- The repair: a check the run can answer
\section{The repair: a check the run can answer}\label{sec:residual}

\begin{planned}{the runtime check, and a run of consecutive lifts}
What a compiler can emit when no verdict fixed before the run will do: a table
with one row per measurement record, computed before the run and looked up
after it, which \cref{fig:two-lane} draws. The rule at one lift, the two
hypotheses it needs, and the closed form for a run of consecutive lifts are
being written; nothing here is claimed.
\end{planned}

\begin{figure}
\centering
\includegraphics{figures/two-lane}
\Description{A two-lane compiler workflow. Before execution the compiler cuts
the program at its lifts, checks the prefix, and emits one proportionality
test per measurement record. When an outcome arrives, execution selects that
record's row and admits or refuses the drop.}
\caption{\revisinginline{The two lanes of \cref{sec:residual}. \textbf{Look at the arrow
between them}: what the compiler emits is a \emph{table}, one row per
measurement record, not a verdict. The left lane is the work
\cref{sec:boundary} leaves open before the run, and it abstains rather than guessing
when the prefix does not leave the temporary constant with a right-invertible
visible action. The right lane needs the run, and it is a lookup.}}
\label{fig:two-lane}
\end{figure}


% Section 9 -- Several temporaries, and lifetimes
\section{Several temporaries, and lifetimes}\label{sec:many}

\begin{planned}{composition, lifetimes, summaries, and borrows}
The criterion is stated for one temporary. What a summary must retain when a
program allocates several---composed into one, tracked a lifetime at a time as
a Rust-style region discipline would~\cite{KlabnikNichols2023Rust}, or crossed
by a procedure boundary---is being written; \cref{sec:eval}'s borrowed-ancilla
row reports only what the checker printed.
\end{planned}


% Section 10 -- Implementation and evaluation
\section{Implementation and evaluation}
\label{sec:eval}

\revising{The preceding sections prove; this one runs what they define. The artifact is one Haskell package with two implementations of one question. \code{DQT} builds the denotation $K$ as a matrix, forms the Kraus
operators, and takes the exact rank of their $d \times m^{2}$ reshaping
(\cref{def:core}), a cost exponential in the width.
\code{S}, the surface language of the corpus, never builds $K$: it runs the
program on $\dim+1$ test inputs, the visible basis plus all-ones, factors
the state at the drop site, and asks whether the resulting lines agree. A
bridge module compares the two, and \textbf{they agree on every catalogue
verdict}; each engine's runtime check is pinned against its own criterion.
They share no code but one exact cyclotomic backend, so the comparison varies
the algorithm and holds the arithmetic fixed. A companion Lean
library over Mathlib mechanises \cref{thm:criterion} from the calculus's
syntax and typing rules, and proves the Gram identity, the temporary's state
at the maximally mixed visible input, that local gates compose and carry a
product input to a product output, and, in abstract form, the constancy
lemma that \cref{sec:boundary} will state. The Haskell suite runs
\revisinginline{58} named statements in the theory layer and 244
checks in the language layer over seven of the table's eight algorithms and
three more (Shor's order finding, Simon's problem, a temporary \textsc{and}
uncomputed by measurement), the $C^{3}X$ borrow being pinned by three checks on its source files.}

\subsection{The catalogue: a producer rule against the criterion}

\revising{Seventeen small programs, each built to make one point and listed in the
supplement, Appendix~B, run against two
rules. \textbf{R1} accepts a drop when every gate is a permutation, the
gate-level reading of ``creates no superposition''. \textbf{R2} adds ``and no
gate writes visible data into the temporary'', the shape of Silq's producer
rule. The criterion accepts \textbf{twelve} of the seventeen. R1 accepts
four, one wrongly: \code{compute-alone} creates no superposition and still
leaks, because it writes visible data into the temporary. R2 is sound and accepts \textbf{two} of the twelve. Deutsch--Jozsa, Bernstein--Vazirani, Grover's oracle and the two-gate program of \cref{sec:intro} are safe and refused by it.}

\subsection{Eight textbook algorithms}

\begin{revisingblock}
We write eight textbook algorithms in \code{S}, compute the verdict at
every $\drop$, and run the runtime check and Silq on each row's smallest
faithful program. Two rows leave the calculus: $C^{3}X$ takes its temporary on
a state it did not prepare, and superdense coding lifts a \emph{visible}
register while the temporary is live. Their verdicts are the checker's,
computed on the same denotation by the same test, and no theorem here covers
them. The Silq column is transcripts against two source builds, \code{a47d974}
and \code{5417dce}; \code{code/silq/TRANSCRIPT.md} says which run used which,
and marks the one attribution it had to infer. The Q\# column is its release
requirement, unenforced on hardware.
\begingroup\centering
\footnotesize
% Narrow columns: ragged right and no hyphenation, so no cell stretches its
% words across the column and no word breaks in half.
\hyphenpenalty=10000 \exhyphenpenalty=10000 \rightskip=0pt plus 1fil
\begin{tabular}{@{}L{.22\linewidth}L{.28\linewidth}L{.20\linewidth}L{.22\linewidth}@{}}
\toprule
algorithm & this theory says & Silq, run & Q\# release \\
\midrule
Deutsch--Jozsa & safe, left in $[|{-}\rangle]$ (\cref{sec:dj})
  & \emph{cannot synthesize forget} & fails at $|{-}\rangle$ \\
Bernstein--Vazirani & safe, $W = X^{s}$
  & the same error & fails at $|{-}\rangle$ \\
Grover oracle, compute--uncompute & safe, left in $[|0\rangle]$
  & accepts, runs & passes \\
$C^3X$, borrowed dirty ancilla & safe on the two loans run, left on the loan's own ray
  & expressible, unchecked & promise, unchecked \\
Phase estimation & safe, left on the eigenstate line
  & the same error & fails at $|u\rangle$ \\
Superdense coding & safe under every record
  & accepts, no drop & passes \\
Teleportation & both drops safe, every record, for the message $|{+}\rangle$
  & accepts, runs & runtime only \\
Repeat-until-success round & \emph{outcome-dependent}, the runtime check
  decides both & accepts, then \emph{bad forget} on $6/20$ seeds & no single
  release ray \\
\bottomrule
\end{tabular}
\par\endgroup

\emph{The disagreements never go the dangerous way, but ``no
disagreement'' is often ``no ruling''.} In these eight rows Silq never accepts an
automatic $\drop$ that this theory refuses, and it refuses three safe ones;
only the Grover row is positive agreement, since superdense coding's encoding has no \code{forget},
the borrowed ancilla is unchecked, and teleportation's survivor is the
receiver's register. Containing a measurement is not the problem:
superdense coding and teleportation lift outcomes and are not
outcome-dependent. On the last row Silq accepts a hand-supplied
value witness and then fails at run time, \emph{bad forget}, on $6$ of $20$
seeds. That refutes the one value we supplied, and we searched for no other;
that no fixed value can serve at all is what \cref{sec:boundary} will show.
Phase estimation is complex, so its verdict is taken up to phase, on the ray
the temporary is left on rather than on a real vector.
These rows measure agreement and coverage, not scale: both procedures are
exponential in register width.

The $C^3X$ row is a program we did not write: Gidney's four-Toffoli
3-controlled $X$ with one \emph{dirty} ancilla, the first case study of Su et al.'s heap logic~\cite{Su2024Heap}. In the surface language:
\begin{prog}[l@{\hspace{1.4em}}l]
vis q1:2, q2:2, q3:2, q:2; newdirty a:2 @ (1,0); & -- the loan: a arrives on |0>\\
gate CCX q3 a q; gate CCX q1 q2 a;                & -- two Toffolis\\
gate CCX q3 a q; gate CCX q1 q2 a; drop a;        & -- again, then return the loan
\end{prog}
Our checker admits the drop, printing the loan's own ray $(1,0)$ and, with
the loan $|{+}\rangle$, that loan's ray. What it does not do is certify the
borrow for \emph{every} lent state: that needs a finite test set, and this
document does not state one.
\end{revisingblock}

% Section 11 -- Related work
\section{Related work}\label{sec:related}

\begin{revisingblock}
\emph{Affine typing and circuit languages.} Selinger and Valiron make
quantum data affine, so discard is a typing restriction rather than a
question~\cite{SelingerValiron2006}, and QWIRE supplies the circuit/host
split~\cite{Paykin2017QWIRE}. Twist and Guppy address the neighbouring
obligation that a value be unentangled from what it does not
own~\cite{Yuan2022Twist,Koch2025Guppy,Koch2025Ownership}. Twist comes
closest: its Corollary~6.2 is a semantic \emph{iff} for purity, weakened in
\S7.6 to a sufficient condition once assertions are added, and its checker
runs a static-plus-runtime split, but it has no completeness theorem and
no proof that the retreat is forced (\cref{thm:criterion} and \cref{sec:boundary}).

\emph{Automatic uncomputation.} Silq~\cite{Bichsel2020Silq} and
Unqomp~\cite{Paradis2021Unqomp} ask how a value was produced
(\cref{sec:intro}). Cavy~\cite{McNally2021Cavy} and
Qurts~\cite{Hirata2025Qurts} track ancilla lifetimes; Cavy states the
obligation semantically---a short-lived ancilla returns
\emph{unentangled}---and pursues it by lifetime-directed uncomputation, its
correctness proof left open.
Qunity~\cite{Voichick2023Qunity}, Reqomp~\cite{Paradis2024Reqomp} and
modular synthesis~\cite{Venev2024Modular} work under error,
space or purity restrictions. ReQWIRE frames the gap: a semantic
characterisation of ancilla validity, then syntactic conditions it calls
sufficient but not necessary~\cite{Rand2019ReQWIRE}.

\emph{The nearest neighbour.} Qutes gives scope-bounded,
liveness-guided uncomputation a semantic account---semantic lifetimes, a
Restoration Invariant, nested scopes---and its boundary is exact where it applies: for a
temporary entangled with a variable that stays live across the boundary,
restoration is possible \emph{if and only if} the subcomputation is
\emph{output-isolable} with respect to that variable, decomposing as $U_L
\circ V_T$ with $V_T$ acting on the temporary
alone~\cite[Def.~5, Lem.~11]{Faro2026Qutes}; a
companion HPDC paper studies the structural parallelism such closed regions
create~\cite{Faro2026Parallelism}. Read literally, Definition~5 puts the two
factors on disjoint registers, so the subcomputation denotes a temporary
factor tensored with a visible one---both gates \emph{local} in the sense of
\cref{sec:calculus}---and that is strictly stronger than drop safety; their worked example
inserts the temporary's own history, reading a live input variable, so we do
not read it as a comparison of strength. Reclamation returns the temporary to
its \emph{initial} state~\cite[Def.~6]{Faro2026Qutes}, and our question is the
other one: may a temporary be dropped \emph{without} restoration?
Deutsch--Jozsa's ancilla ends on $|{-}\rangle$, not where it was allocated,
while \cref{thm:criterion} asks only that it end \emph{somewhere} fixed. Where
a measurement outcome influences what follows their invariant cannot hold past
that point~\cite[Lem.~9]{Faro2026Qutes}; that is \cref{sec:boundary}'s regime,
where the verdict is indexed by the record rather than refused.

\emph{Verification logics.} Q* makes separability an assertion
justifying discard~\cite{Hietala2022Qstar}; Tsai et al.\ compile such
assertions to automata~\cite{Tsai2026Specification}. Su et al.'s quantum BI
logic is sound and relatively complete for dirty-allocation
heaps~\cite{Su2024Heap}. Hoare Meets Heisenberg certifies disposal for a
stabilizer fragment~\cite{Sundaram2026Hoare}. None contradicts
\cref{sec:boundary}: each verifies a given postcondition, our static ceiling;
the impossible boolean is the run-sensitive one, its sound side Chen et
al.'s dead-measurement analysis~\cite{Chen2026DeadMeasurement}.

\emph{Exact conditions at circuit level.} Li et al.\ verify that a
circuit restores its clean or dirty ancillas, by commutation checks on decision diagrams and model counting, and repair faults~\cite{Li2026AncillaSafety}. Liu et al.\ show that deciding whether an uncomputation exists is coNP-complete for reversible Boolean circuits over $\{\mathrm{CCNOT}, X\}$, and coNP-hard wherever those two are polynomial-size constructible~\cite{Liu2026Uncomputation}. Su et al.\ give finite tests for the
identity-on-the-borrowed-qubit condition, two states for classical
circuits~\cite{Su2025Borrowing}, and Bona assumes safe dirty-ancilla
blocks---identity on the borrowed qubit---and schedules semantics-preserving
borrow substitutions~\cite{Xu2026Bona}. All four work on fixed unitary
circuits. A drop asks whether restoration is needed at all, strictly between
restoration and restorability (\cref{sec:intro}) and, with measurement,
dependent on the record (\cref{sec:boundary}).
\Cref{sec:cost} prices the verdict
itself~\cite{Bennett1973Logical,Tanaka2009ExactNonIdentity}.

\emph{The reversible lineage.} In the $\Pi$ family a fractional-type
token reclaims allocated storage~\cite{Chen2020Fractional}, and
de-allocation is a runtime value match~\cite{Chen2021Algebraic}. QuantumPi adds quantum ancillas, hiding and
measurement~\cite{Carette2024QuantumEffect}, but no outcome-indexed
certificate. Our calculus is not $\Pi$: on outcome-dependent commands the
runtime position is forced, and the fractional reuse rule is sound and
strictly weaker (\cref{sec:criterion}). The spooky pebble
game~\cite{Kornerup2025Spooky} measures rather than restores the ancilla,
and its uncomputed ancilla is drop safety at $g = |0\rangle$---restoration
again---with the phase left as the continuation's obligation.
\end{revisingblock}


% Section 12 -- Conclusion, scope, and what is open
\section{Conclusion, scope, and what is open}\label{sec:conclusion}

\revising{A quantum temporary may be dropped when the rest of the program cannot
tell. For a measurement-free typed program three conditions
coincide---weakenable, drop-safe, contextually indistinguishable from a
separable program---and the test is proportionality of the Kraus operators: Kraus rank
one, equivalently a pure temporary at the maximally mixed visible input. Over a gate ring completable in the temporary's dimension it
acquires one side condition, supplied by the Clifford$+T$ ring and its two
complex restrictions and not by its two real ones. Once a command lifts an outcome
the condition survives and the \emph{verdict} does not: one
boolean fixed before the run cannot be right on both records, so a sound one
must refuse a safe run, on exactly the outcome-dependent commands. The missing
argument is the measurement record, so the verdict must become record-aware;
ours is a table the compiler emits, one row per record. Every verdict here is
defined exactly from a selected denotation and computed in exact arithmetic;
\cref{prop:avgcost} alone is a bounded-error promise test. Open: a gate ring
with no column lemma, the positive side of the per-lifetime obstruction, and
whether the runtime check of \cref{sec:residual} is exactly some set of
typing rules. Five boxes marked \emph{Planned} name the
results this draft does not yet carry, each being written.}


% ---------------------------------------------------------------------------
% The bibliography is outside the 25 pages.
% ---------------------------------------------------------------------------
\bibliographystyle{ACM-Reference-Format}
\bibliography{refs}

% No appendix here: appendices are supplemental material, not part of the
% main PDF. They are in supplement.tex, which `make` builds as its own
% document.

\end{document}
